\documentclass[11pt]{article}

\usepackage[T1]{fontenc}
\usepackage{lmodern}
\usepackage[margin=1in]{geometry}

\usepackage{amsmath,amssymb,amsthm}
\usepackage{graphicx}
\graphicspath{{figs/}}
\usepackage{booktabs}
\usepackage{array}
\usepackage[numbers,sort&compress]{natbib}
\usepackage{xcolor}           % before hyperref: blue!50!black is xcolor-only
\usepackage{hyperref}         % load last
\hypersetup{
  colorlinks=true, linkcolor=blue!50!black,
  citecolor=blue!50!black, urlcolor=blue!50!black,
  pdftitle={Accuracy Is Not Auditable: Grounding as the Measurable Property of Language Model Assertion},
  pdfauthor={David Callaghan}
}

\emergencystretch=2em

\renewcommand{\topfraction}{0.9}
\renewcommand{\bottomfraction}{0.8}
\renewcommand{\textfraction}{0.07}
\renewcommand{\floatpagefraction}{0.7}
\setcounter{topnumber}{3}
\setcounter{bottomnumber}{2}
\setcounter{totalnumber}{5}
\usepackage[section]{placeins}

\newcommand{\Out}{\mathrm{O}}          % O_S(phi)      : system S asserts phi
\newcommand{\Grd}{\mathrm{G}}          % G_S(phi | c)  : phi is grounded at c

\newcommand{\Hall}{\operatorname{Hallucination}}
\newcommand{\Err}{\operatorname{Error}}
\newcommand{\Misapplied}{\operatorname{Misapplied}}
\newcommand{\Abstain}{\operatorname{Abstain}}

\newcommand{\Supports}{\operatorname{Supports}}
\newcommand{\Applicable}{\operatorname{Applicable}}
\newcommand{\Current}{\operatorname{Current}}
\newcommand{\Authentic}{\operatorname{Authentic}}
\newcommand{\Traceable}{\operatorname{Traceable}}
\newcommand{\Resolved}{\operatorname{Resolved}}
\newcommand{\TrainedOn}{\operatorname{TrainedOn}}

\newcommand{\nentails}{\;\not\Rightarrow\;}

\newcommand{\Grounded}{\operatorname{Grounded}}
\newcommand{\Versioned}{\operatorname{Versioned}}
\newcommand{\Authoritative}{\operatorname{Authoritative}}
\newcommand{\Entails}{\operatorname{Entails}}
\newcommand{\Sensitive}{\operatorname{Sensitive}}

\newcommand{\Warranted}{\operatorname{Warranted}}
\newcommand{\Verified}{\operatorname{Verified}}
\newcommand{\Authorized}{\operatorname{Authorized}}
\newcommand{\True}{\operatorname{True}}

\newcommand{\Sufficient}{\operatorname{Sufficient}}
\newcommand{\NoDefeater}{\operatorname{NoDefeater}}
\newcommand{\Contested}{\operatorname{Contested}}

\newcommand{\standing}{\alpha}
\newcommand{\DefeaterSearch}{\mathcal{D}}

\theoremstyle{plain}
\newtheorem{proposition}{Proposition}
\newtheorem{corollary}{Corollary}
\theoremstyle{definition}
\newtheorem{definition}{Definition}
\newtheorem{fact}{Fact}

\title{Accuracy Is Not Auditable:\\
       Grounding as the Measurable Property of\\
       Language Model Assertion}
\author{%
  David Callaghan\\
  \texttt{mr.david.callaghan@gmail.com}
}
\date{\today}


\begin{document}
\maketitle

\begin{abstract}
Language model deployments in regulated environments are commonly validated by
comparing outputs against reference answers. This paper argues that the practice cannot
deliver what sign-off requires, and identifies the property that can.

The first of two connected claims is a redefinition of hallucination, indexed to a
decision context. Existing definitions ask whether an assertion is supported by a
supplied source; this one asks the prior question of whether that source governs the
case, by requiring support to be applicable at a resolved context of jurisdiction, valid
time, material facts, and intended use, and by requiring the assertion to have been
produced because of the evidence rather than merely to agree with it. The consequence is
a failure class---authentic, authoritative, correctly cited evidence that does not govern
the decision---which faithfulness and attribution metrics score as clean by construction,
since the entailment genuinely holds. Unlike the truth-conditional definition, this one
can be evaluated at the point of assertion without a reference answer, and so functions as a
gate rather than only as a score.

The second claim is that the same architecture makes accuracy uncertifiable. The system
computes a distribution over token sequences conditioned on token sequences and applies
no canonical form, so a correct answer observed once carries no guarantee that a
materially equivalent request will be answered the same way. Accuracy is a property of a
sample and the sample does not generalise. I prove that accuracy places no bound whatever
on the grounding rate---the two vary independently over the whole unit square---so no
amount of reference-based evaluation constrains how often a system asserts beyond its
evidence.

The constructive half derives a measurement protocol, nineteen system invariants, a
reference architecture for governed retrieval, and an account of why the division of
ownership across data, policy, platform, validation, and records functions is itself a
principal source of failure.
\end{abstract}


\section{Introduction}
\label{sec:intro}

Consider a decision you have to be able to defend two years from now.

A clinical operations team asks a language model system whether a compound is
approved for a paediatric indication. The system answers. Someone acts on the
answer. Later---in an audit, an inspection, a dispute---the question arrives:
what was that answer based on? Not whether the model is generally accurate. Which
document, in which version, as of which date, under which jurisdiction, supported
this specific claim.

Most current validation practice cannot answer that question, and the reason is
not that the practice is careless. It is that the practice measures the wrong
thing. The dominant approach evaluates a model by sampling questions with known
answers and scoring the outputs against them---\emph{reference-based evaluation},
in the terms of the natural-language-generation literature. It produces a
number---accuracy, or its complement, often reported as a hallucination
rate---and that number is then used as evidence of fitness for deployment.

The central argument of this paper can be stated in one sentence. \textbf{In
regulated settings, the relevant question is not merely whether an answer is
accurate. It is whether the organisation can reconstruct why this particular
assertion was produced, identify the authoritative and versioned evidence on which
it rests, establish that the evidence applies to the resolved case context, verify
the asserted relation, and show that reliance was authorised under the institution's
governing rules.} Accuracy is one input to that question and settles none of it.

This paper argues that reference-based evaluation cannot support the use it is put
to,
for two reasons that compound.

The first is a classification problem. Pooling every wrong answer into one
metric conflates failures with different causes and different remedies. A system
that faithfully reports an erroneous source has a data problem. A system that
fabricates a plausible-looking figure has a generation problem. A system that
correctly reports a regulation that was superseded last year, or that applies in
a different jurisdiction, has a context-resolution problem. These are three
separate defects, addressed by three separate teams, on three separate budgets.
One number does not tell you which one to spend.

The second reason is worse, and it is the argument of this paper.
Section~\ref{sec:canonical-form} shows that the system applies no canonical form
to its inputs. It computes a conditional distribution over token sequences given
token sequences; its inputs are strings; and two semantically identical
queries---a contraposed conditional, a quantity written in digits rather than
words, an active rather than passive phrasing---are two distinct points in input
space. Nothing in the architecture guarantees they map to the same output, and
measurement confirms they frequently do not. A database normalises before it
compares. A compiler canonicalises before it optimises. This system does
neither.

The consequence for validation is direct. If a test case passes, you have
learned that one string produced an acceptable output on one occasion. You have
not learned that the semantically identical string will, or that the same string
will next week, or that the same string will later in the same conversation. A
passing test does not generalise to the paraphrase, so the test suite does not
certify the behaviour it appears to certify. Accuracy is a property of a sample,
and this sample does not support the inference that validation requires of it.

That is a negative result, and on its own it would be unhelpful. The
constructive claim is that a different property does not have this defect.
\emph{Grounding}---whether an assertion is supported by evidence applicable to
the decision at hand, through a mechanism actually sensitive to that
evidence---is a property of the relation between an assertion and an evidential
base, not of a sampled output. It is stable under paraphrase in the way accuracy
is not, because it is a fact about the pipeline rather than about a draw from it.
It is therefore the property a governance regime can audit.

Section~\ref{sec:auditable} makes this precise, and proves the sharpest form of
the point: accuracy and the grounding rate vary independently over the entire
unit square. No accuracy figure constrains the rate at which a system asserts
beyond its evidence, in either direction. A system can be highly accurate and
almost entirely ungrounded, and the second fact is the one that predicts what
happens when the questions change.

\paragraph{What this paper is not arguing.} Not that these systems are
unreliable in a way that precludes use. I build them, in exactly the regulated
settings this paper is about, and they work. Not that accuracy is uninteresting;
it is interesting, and it is the wrong thing to certify against. And not that
better models will not help. They will, on some axes and not others, and
Section~\ref{sec:architecture} says which.

\paragraph{Contributions.} Six, in the order they appear.

\begin{enumerate}
  \item \textbf{A redefinition of hallucination, indexed to a decision context.}
        Existing definitions ask whether an assertion is supported by a supplied
        source. This one asks the prior question---whether that source governs the
        case---by requiring support to be applicable at a resolved context of
        jurisdiction, valid time, material facts, and intended use, and by
        requiring the assertion to have been produced \emph{because of} the evidence
        rather than merely to agree with it. The consequence is a failure class that
        faithfulness and attribution metrics score as clean by construction, and a
        definition evaluable at the point of assertion without a reference answer
        (Sections~\ref{sec:taxonomy}, \ref{sec:novelty}).
  \item A five-category taxonomy following from it---\textsc{Sound},
        \textsc{Lucky}, \textsc{Error}, \textsc{Fabrication}, and
        \textsc{Misapplied}---in which the fifth, assertion grounded in authentic
        but inapplicable evidence, is a distinct control failure with its own
        remedy and owner (Sections~\ref{sec:taxonomy}, \ref{sec:cell-v}).
  \item A proof that aggregate accuracy neither identifies nor bounds the grounding
        rate, so a validation programme reporting accuracy alone cannot distinguish
        radically different reliance profiles
        (Proposition~\ref{prop:independence}).
  \item An architectural account of why the failures are structural, grounded in the
        specific architecture and training regime these systems inherit, and an
        argument that the absence of a canonical form makes accuracy uncertifiable
        rather than merely incomplete
        (Sections~\ref{sec:architecture}, \ref{sec:canonical-form}).
  \item A set of nineteen invariants separating what assertion behaviour does not
        supply from what a system must be built to supply, each tied to a control
        (Section~\ref{sec:invariants}).
  \item A reference architecture for governed retrieval, a definition of
        institutional warrant, and an account of why the division of ownership
        across data, policy, and platform functions is itself a principal cause of
        failure (Sections~\ref{sec:governed-rag}--\ref{sec:three-teams}).
\end{enumerate}

\paragraph{Scope.} This paper is about one class of system: models that emit
natural-language assertions by autoregressive sampling over a token vocabulary,
fitted by next-token likelihood on large text corpora, and optionally post-trained
on human preferences. That is the decoder-only transformer language model
established at scale by the GPT-3 generation and everything built on that pattern
since. It is not about deep learning generally. Generative adversarial networks and
diffusion models are out of scope, as are encoder-only classifiers, rankers, and
extractors, and the conventional predictive models that make up most of an
enterprise analytics estate. Those components assert nothing, so grounding is not a
coherent question about them---which is not to say they are easy to validate;
Section~\ref{sec:not-easier} is explicit that they share several of the problems
described here. Section~\ref{sec:scope} states the boundary
precisely and says why it falls where it does.

\paragraph{Structure.} Section~\ref{sec:premise} states the premise the argument
rests on. Section~\ref{sec:architecture} identifies the systems in scope and
derives, from their architecture and training regime, why the failures this paper
describes are structural. Section~\ref{sec:taxonomy} redefines
hallucination as ungrounded assertion and gives the five categories that follow.
Section~\ref{sec:independence} proves that accuracy and the grounding rate are
independent.
Section~\ref{sec:grounding-context} defines the resolved context and the six
dimensions of grounding, and explains why retrieval and citation do not supply
them. Section~\ref{sec:canonical-form} develops the canonical-form argument and its
consequence for validation. Section~\ref{sec:invariants} states the invariants.
Section~\ref{sec:auditable} sets out why grounding rather than accuracy is
auditable, and where it can be measured. Sections~\ref{sec:governed-rag}
to~\ref{sec:three-teams} give the constructive architecture, the definition of
warrant with reliance tiers, and the organisational analysis.
Section~\ref{sec:evaluation} gives an evaluation protocol and
Section~\ref{sec:related} positions the work.

\paragraph{Notation.} Formal notation is used where it makes a claim checkable
and avoided where it would only make one look rigorous. Two operators, five
defined categories, one proposition with a two-line proof. Every definition is
stated over observable behaviour, so that any claim made here can be tested
against a running system.


\section{Premise: What Generation Does Not Establish}
\label{sec:premise}

One premise carries the rest of the paper, and it is deliberately modest. A
system's producing a statement is a different thing from that statement being
supported by evidence, applicable to the case at hand, checked, and permitted for
use. Any architect will grant the distinction in every other context: a report that
renders is not a report that reconciles, and a value that populates a field is not
a value that passed validation. The premise says only that generated text is
subject to the same distinction, and that the generation event is not what closes
it.

The argument below needs no position on what a model understands or represents
internally. It needs the distinction stated formally, which is the business of this
section.

\subsection{Notation}
\label{sec:notation}

\begin{center}
\begin{tabular}{@{}ll@{}}
\toprule
$M$ & the language model: a parameterised generative function \\
$S$ & the deployed system: $M$ together with its decoding, retrieval, policy,
      and control layers \\
$x$ & an input---a prompt, a structured request, a conversation state \\
$p, q$ & propositions or claims \\
$c$ & the resolved context of a decision \\
$E$ & an evidentiary basis \\
$V$ & a verification procedure \\
$I$ & an institution or accountable organisation \\
\bottomrule
\end{tabular}
\end{center}

The distinction between $M$ and $S$ is load-bearing rather than pedantic. Every
property this paper identifies as necessary for reliance is a property of $S$, and
none is a property of $M$. A model has no version of an audit record, an
authorization, or a resolved context; a system can have all three. Where the
literature reports a result about models, I say $M$; where the object of
governance is at issue, I say $S$.

\subsection{The resolved context}
\label{sec:resolved-context}

A great deal of enterprise information is true only relative to a governing frame.
Write
\begin{equation}
  c = \{\, j,\ t_v,\ f,\ u \,\},
  \label{eq:context}
\end{equation}
where

\begin{center}
\begin{tabular}{@{}l>{\raggedright\arraybackslash}p{10.6cm}@{}}
\toprule
$j$ & jurisdiction or governing regime: the legal, regulatory, or policy authority
      whose rules apply \\
$t_v$ & \emph{valid time}: the effective date the assertion is about---the as-of date
      against which the governing rules are to be assessed \\
$f$ & material case facts: patient facts, customer facts, transaction facts,
      product facts, or whatever else is decision-relevant \\
$u$ & intended use: the decision class, the reliance context, what the answer will
      be used to do \\
\bottomrule
\end{tabular}
\end{center}

A second time dimension is required and it does not belong in $c$. Write $t_k$ for
\emph{transaction time}: the state of the evidential base the system actually had
access to---which corpus release, ingested to what date. The distinction is the
ordinary bitemporal one and it is not decoration.

$t_v$ is a property of the \emph{question}: which date's rules govern this decision.
$t_k$ is a property of the \emph{answer's provenance}: what the system could have
known when it answered. Correctness is assessed against $t_v$; reconstruction requires $t_k$; and their
relation is checkable in \emph{both} directions, which is the whole reason two
timestamps are needed:
\begin{align}
  t_k \prec t_v &\;\Longrightarrow\; \text{the base may omit instruments in force at }
    t_v \quad \text{(omission)}, \label{eq:bitemporal-stale}\\
  t_k \succ t_v &\;\Longrightarrow\; \text{the base may contain instruments not in
    force at } t_v \quad \text{(anachronism)}. \label{eq:bitemporal-anachronism}
\end{align}

Omission is the familiar case: a rule amended before $t_v$ but ingested after $t_k$
leaves an assertion ungrounded in a way no inspection of the retrieved evidence will
reveal, because the missing instrument is missing.

Anachronism is the other half and in regulated work it is the more dangerous of the two,
because the answer is confidently wrong rather than incompletely right. A corpus current
to today, answering a question as-of two years ago, will retrieve provisions that had
not yet been made. Applying a rule to a transaction that preceded it is a familiar and
serious error in its own right, and the system will produce a clean citation to a real
instrument while committing it.

Neither failure is prevented by having a recent corpus, and both are prevented by the
same condition. $\Applicable(E, c)$ must be read to require \emph{interval containment}
on the temporal component of $c$: the instrument's own validity period includes $t_v$.
That excludes the not-yet-enacted and the already-superseded together, and it is a
filter on the evidence rather than a property of the corpus. Temporal validity is
therefore part of applicability rather than a dimension of its own, which is why the
seven conditions of Definition~\ref{def:grounded-evidence} do not list it separately. Both directions are also detectable at the aggregate level
without knowing what any amendment said: compare the corpus ingest watermark against
the decision date, per source class, and flag either sign of the difference. Systems that record only one timestamp cannot
perform the comparison, and most record only one.

$\Resolved(c)$ holds when the components of \eqref{eq:context} are explicit rather
than assumed or inferred without validation. $t_k$ is not resolved by the requester;
it is recorded by the system, and it belongs to I9 and I14 rather than to I10. The distinction between explicit and inferred is the whole of
it: a system that guesses the jurisdiction correctly has not resolved the context,
because nothing in the record establishes that the guess was checked, and an
auditor cannot distinguish a correct guess from a lucky one.

The truth predicate is correspondingly indexed. $\True(p \mid c)$ is truth relative
to a resolved context, and the relation that fails is:
\begin{equation}
  \True(p \mid j_1, t_1, f_1, u_1)
  \nentails
  \True(p \mid j_2, t_2, f_2, u_2).
  \label{eq:no-transfer}
\end{equation}

\subsection{The non-entailments}

Write $\Pr_M(\mathrm{utter}(p) \mid x)$ for the probability that $M$, given input
$x$, produces an assertion with content $p$. The premise is that a high value of
this quantity does not establish any of the properties on which reliance depends:

\begin{align}
  \Pr_M(\mathrm{utter}(p) \mid x) \text{ high}
    &\nentails \True(p \mid c),            \label{eq:ne-true}\\
  \Pr_M(\mathrm{utter}(p) \mid x) \text{ high}
    &\nentails \Supports(E, p, c),         \label{eq:ne-supported}\\
  \Pr_M(\mathrm{utter}(p) \mid x) \text{ high}
    &\nentails \Applicable(E, c),          \label{eq:ne-applicable}\\
  \Pr_M(\mathrm{utter}(p) \mid x) \text{ high}
    &\nentails \Verified(V, p, c),         \label{eq:ne-verified}\\
  \Pr_M(\mathrm{utter}(p) \mid x) \text{ high}
    &\nentails \Warranted_I(p \mid c).     \label{eq:ne-warranted}
\end{align}

Compactly, and in the order in which an institution needs them:
\begin{equation}
  \Pr_M(\mathrm{utter}(p) \mid x) \text{ high}
  \nentails \True(p \mid c)
  \nentails \Supports(E, p, c)
  \nentails \Warranted_I(p \mid c).
  \label{eq:ne-chain}
\end{equation}

These are non-entailments, and the distinction matters because the claim is
frequently misread as stronger than it is. They do not assert that the properties
fail to co-occur. A model can and routinely does produce propositions that are
true, well supported, applicable, verifiable, and fit for reliance. The claim is
that the generation event is not what establishes those properties, so the
generation event cannot be what an institution points to when asked to defend
them. Each arrow in \eqref{eq:ne-chain} marks a gap that something other than the
model has to close.

Nor is the claim that the properties are unrelated. High generation probability is
correlated with truth on the distribution the model was fitted to, sometimes
strongly. Correlation on a training distribution is not entailment, and
Sections~\ref{sec:architecture} and~\ref{sec:canonical-form} concern what happens
when the distribution of interest is not that one.

\subsection{What follows}

The premise is weak, and that is its value: it can be granted without conceding
anything contentious, and it is sufficient for everything in this paper.

Two clarifications, because both are easy to misread from the non-entailments
alone. The claim is not that model output is unreliable. Output is often excellent,
and the argument concerns what can be established about an individual assertion
rather than what can be expected of a population of them. And the claim is not that
the properties in \eqref{eq:ne-chain} are rare. In a well-built system most
assertions have all of them. The point is that they are established elsewhere, by
components that can be pointed to.

That is the whole of the constructive programme. If generation does not establish
support, applicability, verification, or authorization, then those properties are
supplied by something in $S$ that is not $M$---and identifying, building, and
measuring that something is what the rest of this paper is about.


\section{The Systems in Scope}
\label{sec:architecture}

Before classifying what these systems get wrong, it is worth being exact about
which systems are meant. The argument that follows depends on specific properties
of one architecture and one training regime, and it does not generalise to deep
neural networks at large.

\subsection{What is in scope, and what is not}
\label{sec:scope}

In scope are systems with all four of the following properties.

\begin{enumerate}
  \item They emit natural-language assertions---propositional content a reader can
        act on---as their output.
  \item They produce that output by autoregressive sampling over a fixed token
        vocabulary, one token at a time, each conditioned on the preceding ones.
  \item Their parameters are fitted by next-token likelihood on large text
        corpora.
  \item They are optionally post-trained on human demonstrations and preference
        comparisons.
\end{enumerate}

Concretely: the decoder-only transformer language models established at scale by
the GPT-3 generation \citep{brown2020}, extended by instruction tuning and
preference optimisation \citep{ouyang2022}, and everything built on that pattern
since. That is the point at which these systems became plausible components of a
business process, and it is the regime current deployments inherit.

Out of scope, and worth naming because ``deep learning'' is often treated as one
thing:

\begin{itemize}
  \item \textbf{Generative adversarial networks} and \textbf{diffusion models}.
        Both are generative, and neither generates by autoregressive token
        sampling---an adversarial objective and an iterative denoising objective
        respectively---and neither emits propositional content that an
        organisation relies on as an assertion. A synthesised image makes no claim
        that could be grounded or ungrounded.
  \item \textbf{Encoder-only and discriminative models.} A masked language model
        used as a classifier, a ranker, or an extractor produces a label or a
        score, not an assertion. Grounding is not a coherent question about a
        label: a class prediction does not claim anything that evidence could
        support or fail to support.
  \item \textbf{Conventional predictive and statistical models.} Regression,
        gradient-boosted trees, and the rest of the enterprise analytics estate
        are out of scope for the same reason: they do not assert.
\end{itemize}

\subsection{What the system computes}
\label{sec:what-computed}

The computation has four stages, and the first is where the argument of
Section~\ref{sec:canonical-form} originates.

\paragraph{Tokenisation.} Text is segmented into tokens drawn from a fixed
vocabulary $\mathcal{V}$, typically by byte-pair encoding. This is a purely
lexical operation. It is worth pausing on what it means for the rest of the stack:
\texttt{10}, \texttt{ten}, and \texttt{10.0} are distinct token sequences, as are
\texttt{Q1 2024} and \texttt{first quarter of 2024}. No semantic normalisation
occurs at this boundary, or at any later one.

\paragraph{Embedding.} Token identifiers are mapped to vectors by an embedding
matrix $W_E \in \mathbb{R}^{|\mathcal{V}| \times d}$, and position information is
added. An input of $n$ tokens becomes a matrix in $\mathbb{R}^{n \times d}$.

\paragraph{Causal self-attention.} Each of $L$ layers computes, per head,
\begin{equation}
  \operatorname{Attention}(Q, K, V)
  = \operatorname{softmax}\!\left(\frac{QK^{\top}}{\sqrt{d_k}} + \Lambda\right) V,
  \qquad
  \Lambda_{ij} = \begin{cases} 0 & i \ge j,\\ -\infty & i < j,\end{cases}
  \label{eq:attention}
\end{equation}
where $Q, K, V$ are linear projections of the layer input, $d_k$ scales the logits
to keep the softmax out of saturation, and $\Lambda$ is the causal mask that
prevents position $i$ from attending to positions after it
\citep{vaswani2017}. Attention output is combined with a position-wise
feed-forward transformation, with residual connections and normalisation.

\paragraph{Unembedding and factorisation.} The final layer is projected to logits
over $\mathcal{V}$ and normalised, giving $P(x_k \mid x_{<k}; \Theta)$. The causal
mask is what licenses reading the whole as an autoregressive factorisation:
\begin{equation}
  P(x_1, \ldots, x_n; \Theta) = \prod_{k=1}^{n} P(x_k \mid x_{<k}; \Theta).
  \label{eq:factorisation}
\end{equation}

Now state the signature of the resulting map plainly, because everything in
Section~\ref{sec:canonical-form} follows from it:
\begin{equation}
  S : \mathcal{V}^{*} \longrightarrow \Delta(\mathcal{V}^{*}),
  \label{eq:signature}
\end{equation}
a function from token sequences to distributions over token sequences. The domain
is $\mathcal{V}^{*}$. Not propositions, not normalised queries, not parsed
intents, not records with identity---token sequences. There is no slot in the
parameterisation for a proposition, and no stage between the tokeniser and the
unembedding at which two encodings of the same claim are brought to a common form.
This is not an omission that could be patched at the margins; it is what the
architecture is.

There is direct evidence that content is stored in a form tied to its surface
presentation rather than propositionally. Models trained on statements of the form
``$A$ is $B$'' assert them reliably and fail to answer a query directed at $B$, even
though the fact is the same fact \citep{berglund2024}. The result is about what was
learned rather than about how two phrasings of one request are treated---a
distinction Section~\ref{sec:canonical-form} keeps, and this is not an instance of
the invariance failure discussed there. What it shows is more fundamental and more
useful here: if a fact were represented as a proposition, its availability would not
depend on the order in which it was encountered. That it does is evidence there is no
propositional representation to depend on.

\subsection{The pretraining objective}
\label{sec:pretraining}

Parameters are fitted by minimising cross-entropy on next-token prediction over a
large corpus:
\begin{equation}
  \mathcal{L}_{\mathrm{pre}}
  = -\sum_{k} \log P\!\left(x_k \mid x_{<k}; \Theta\right).
  \label{eq:objective}
\end{equation}
For the GPT-3 generation this was applied at a scale of order $10^{11}$ dense
parameters over a corpus of order $10^{11}$ tokens drawn from filtered web crawl,
curated web text, books, and encyclopaedic sources \citep{brown2020}.

Read \eqref{eq:objective} as a specification and audit it for what is absent. The
claim to make here is narrower than it is tempting to make. Fitting human text does
make the model sensitive to truth-correlated structure---that is why these systems are
useful---so it is wrong to say the objective is blind to truth. What
\eqref{eq:objective} contains is no term that \emph{rewards truth over plausible
falsehood where the two diverge}. On the training distribution they mostly coincide
and the loss cannot tell them apart; off it they separate and the loss remains
indifferent. Beyond that, it contains no term that is a function of evidence. No term that is a function of
evidence, or of the provenance of the text being predicted. No term that is a
function of the context in which an assertion might later be relied upon. And no
term that is minimised by declining to produce a continuation---the objective is
defined over every position, and an abstention is simply a different continuation
with its own probability.

What \eqref{eq:objective} does reward, and rewards very effectively, is assigning
high probability to continuations resembling the training distribution. A fluent
falsehood in a familiar genre resembles the training distribution closely. This is
not a criticism of the objective, which does what it was designed to do; it is an
observation about scope. The system was fitted for \emph{plausible continuation},
and plausible continuation and warranted assertion are different targets that
coincide on the training distribution and separate off it.

One consequence is worth drawing now, because it is the origin of two of the
five failure categories Section~\ref{sec:taxonomy} sets out. Off the training
distribution---an entity never described, a figure never published, a combination
of jurisdiction and effective date that never co-occurred in the corpus---the
system must still emit tokens. It has no representation of the fact that it is now
extrapolating rather than reporting, because \eqref{eq:objective} gave it no
reason to construct one. It produces the most probable continuation, which is
fluent and which rests on nothing. Whether such a continuation happens to be true
is a separate matter, and Section~\ref{sec:taxonomy} argues that the distinction
between those two outcomes is far less important than it looks.

\subsection{In-context learning, and why the prompt is the whole control surface}
\label{sec:icl}

The property of this generation of models that made enterprise use plausible was
not raw perplexity. It was in-context learning: task behaviour induced by
examples and instructions placed in the prompt, with no gradient update
\citep{brown2020}. That is what made a general-purpose model deployable against a
specific business task without a training programme per task, and it is why these
systems entered production at all.

It also has a consequence that is easy to miss and directly relevant to
governance. If task behaviour is induced through the prompt, then the prompt is the
control surface---and by \eqref{eq:signature} the prompt is a token sequence. All
conditioning, including everything an architect would think of as configuration,
policy, or scope, enters the computation as surface form and is processed by the
same mechanism as the content. There is no separate typed channel for constraints.

This is the architectural root of the failure examined in
Section~\ref{sec:canonical-form}. Other channels exist---system prompts,
fine-tuning, adapters, logit biases, tool schemas---but the per-request control
surface, the one that varies with what a user asks, is the prompt. When that surface
is a string, and
the map's domain is strings, sensitivity to the form of the string is not a defect
layered on top of the design. It is the design operating as specified.

\subsection{Instruction tuning and preference optimisation}
\label{sec:rlhf}

Deployed systems add post-training. Supervised fine-tuning on demonstrations is
followed, in the regime established by \citet{ouyang2022}, by preference
optimisation: a reward model $r_\phi$ is fitted to human comparisons between
candidate outputs, and the policy is optimised against $r_\phi$ with a divergence
penalty against the supervised reference,
\begin{equation}
  \max_{\Theta}\ \mathbb{E}\bigl[r_\phi(x, y)\bigr]
  - \beta\, \mathbb{D}_{\mathrm{KL}}\!\left(\pi_\Theta(y \mid x)
    \,\|\, \pi_{\mathrm{ref}}(y \mid x)\right).
  \label{eq:rlhf}
\end{equation}

This changes the output distribution substantially, and it is why current systems
are usable where their pretrained predecessors were not. For the purposes of this
paper, three properties of \eqref{eq:rlhf} matter.

\paragraph{The reward is a function of the output, not of the evidence.}
$r_\phi(x, y)$ is fitted to preferences over $(x, y)$ pairs. Nothing in its domain
is the evidential base, the source version, or the resolved context. An assertion
and its ungrounded twin---identical text, different provenance---are
indistinguishable to the reward model, and therefore receive identical
optimisation pressure.

The precise limitation is worth stating carefully, because the strong version is
false. Preference optimisation \emph{can} select for the observable correlates of
grounding: raters shown a well-cited answer and a bare one will prefer the former,
and instructions to raters can make citation quality part of the comparison. What
$r_\phi$ cannot do is \emph{verify} the relation, because $E$ is not in its domain.
It can learn that citation-shaped output is preferred; it cannot learn that a
citation supports the claim attached to it, except through whatever raters happened
to check at labelling time---which does not generalise to evidence they never saw.
So the pressure is toward the appearance of grounding, and the gap between appearance
and relation is exactly where an auditor has to look.

\paragraph{The pressure is toward assertion.} Human raters comparing outputs
prefer answers that are fluent, confident, complete, and helpful. A hedged answer,
a request for missing context, or a refusal reads as less helpful than a direct
one, and will usually lose the comparison unless raters are specifically
instructed otherwise. The gradient therefore points away from abstention. The
sycophancy results discussed in Section~\ref{sec:canonical-form} are a measured
consequence of this structure \citep{sharma2024}, not an incidental defect.

\paragraph{The result is better-calibrated form, not better-grounded content.}
Post-training teaches the surface conventions of careful assertion---citation
shapes, hedging vocabulary, structured justifications---because those conventions
are present in preferred outputs. Teaching the relation those conventions signal is
harder and is bounded by what raters could verify: preference data can carry a signal
about whether citations checked out on the examples that were labelled, and cannot
carry one about evidence nobody looked at. So the conventions generalise and the
relation does not. For an auditor this is the least comfortable finding in the section,
because the direction of the gap is unfavourable: post-training reliably improves how
grounded an output \emph{looks} and improves the grounding rate of \eqref{eq:rate}
only incidentally, which widens rather than narrows the distance between apparent and
actual auditability.

\subsection{Three components that are absent}
\label{sec:missing}

Stated as a gap analysis, three things a warranted-assertion system would contain
are absent from the forward pass, and each maps to a category. These are claims about
$M$, not about $S$. A deployed system may well have an evidence gate, a context
resolution step, and an abstention branch---indeed Section~\ref{sec:invariants}
requires all three, as I11, I10, and I17. The point is that none of them comes from the
model, so each must be built, and a system assembled without them inherits the gaps
below by default rather than by choice.

\paragraph{No evidence gate.} Nothing in the computation consults an evidential
base, forms a judgement about whether it supports the candidate claim, and
conditions emission on the answer. Retrieval places evidence \emph{in the
context}, which raises the probability of continuations consistent with it, but
that is influence, not gating: by \eqref{eq:signature} the retrieved text is more
tokens, conditioning a distribution, not a precondition on output. This is why ungrounded
assertion remains available in retrieval-augmented systems, and why the
sensitivity condition of Definition~\ref{def:grounding} is a genuine constraint
rather than a description of existing behaviour.

\paragraph{No context resolution.} Nothing establishes $c$ before answering.
``Is this approved?'' is syntactically complete while leaving jurisdiction,
effective date, material facts, and intended use unspecified, and there is no
mechanism whose function is to notice the omission. This is the failure
Section~\ref{sec:grounding-context} develops as misapplied evidence.

\paragraph{No abstention primitive.} There is no state corresponding to ``I have
not established this.'' Abstention, where it occurs, is a generated
continuation---the tokens of a refusal, produced because refusal-shaped text was
probable in that context, not because a support check failed. This is why
abstention behaviour is sensitive to phrasing: it is an output, not a branch. Any
system that must reliably decline needs the decision to sit outside $M$.

\subsection{A floor, not merely a gap}
\label{sec:floor}

It might be hoped that the ungrounded categories shrink toward zero with better
training. There is a formal reason to expect otherwise.
% TODO(verify): check this gloss against the statement of the theorem before
\citet{kalai2024} show that a calibrated language model must produce incorrect
assertions at a positive rate on facts appearing rarely in its training
distribution: a model well calibrated on the frequency of rare facts is thereby
committed to generating some of them wrongly, and a model that avoided this would
be miscalibrated in a way that costs it elsewhere.

The engineering reading is that the ungrounded cells have a floor above zero that
is not a defect of any particular system, and that the floor is highest exactly
where enterprise questions live---specific entities, specific dates, specific
jurisdictions, all of them rare in a general corpus. If the generation path cannot
be made not to produce ungrounded assertions, then the surrounding system must be
able to \emph{detect} them. That is the argument for measuring grounding directly
rather than waiting for the model to improve.

\subsection{What this predicts about better models}
\label{sec:better-models}

The analysis supports specific predictions, which matter for roadmap decisions
because they are not the predictions usually made.

Scale and post-training reliably reduce outright fabrication on questions
resembling the training distribution. They do not address a corpus that is wrong,
which is a property of the data rather than the model. They do not address
inapplicable evidence with the context it lacks. Capability gains can make a model
more likely to \emph{ask} for a missing jurisdiction or effective date, which helps;
they cannot make asking reliable, and they cannot supply the answer, which has to come
from the request or from a system that resolves it. And their effect on lucky
guessing runs in the wrong direction: a better model guesses better, which raises
accuracy while leaving the assertion exactly as disconnected from evidence as
before. Section~\ref{sec:taxonomy} gives these four outcomes names and
Section~\ref{sec:remedies} assigns them owners.

One concession is owed here, and it is specific. A system that samples several
candidate outputs and selects among them with a verifier does have a check outside the
forward pass, and where the verifier is sound---a proof checker, a test suite, a
constraint solver---that check is a genuine correctness oracle. Such a system satisfies
something structurally similar to the external gate this paper argues for in
Section~\ref{sec:constrained-output}, and it would be wrong to claim the mechanism is
absent from systems that plainly have it.

What does not transfer is the availability of the oracle. Verification of that kind
requires correctness to be cheaply decidable against a formal specification, which is
why it works for mathematics, code, and constraint satisfaction. There is no oracle
for whether a particular regulatory instrument is the governing one for this
jurisdiction at this effective date, applied to this population, for this decision
class. The question is not undecidable in some deep sense; it is simply not settled by
computation over the assertion. It is settled by evidence and by authority, which is
what Sections~\ref{sec:grounding-context} and~\ref{sec:governed-rag} are about. So
outcome verification is a real mechanism that addresses a class of problems this paper
is not about, and its success there is not evidence about the class it is.

Longer context windows, retrieval, and tool use change what the system can do without
changing \eqref{eq:signature}. The domain
remains $\mathcal{V}^{*}$; there is still no canonicalisation stage; the reward
still cannot see provenance. Retrieval is the most consequential of them for this
paper's purposes, and its contribution is not primarily accuracy---it is that an
identified evidential base makes the grounding rate observable at all, as
Section~\ref{sec:measurability} explains.

\subsection{How this compares with other model families}
\label{sec:not-easier}

Section~\ref{sec:scope} excluded discriminative and predictive models on the ground
that grounding is not a coherent question about a label. With the architecture now in
view the comparison can be made properly, and it should be, because the exclusion is
about \emph{whether grounding applies} and not about whether those systems are easy to
validate. It would be wrong to suggest that discriminative
and predictive models are unproblematic, and the overlap with this paper's categories
is worth naming, because it locates the argument honestly.

A classifier that has learned a spurious feature---the hospital identifier rather
than the pathology, the formatting artefact rather than the content---is right for
the wrong reason. That is the same defect this paper calls a \emph{lucky} assertion:
correct output, no supporting relation, and no reason to expect the next case to go
the same way. A classifier whose label flips under an imperceptible perturbation has
failed a form-invariance requirement. A classifier degrading under distribution
shift has exactly the problem of a sample that does not generalise. And per-instance
attribution is weak for most model families, which is a traceability failure.
Enterprise ML monitoring exists because none of this is solved.

What is different is narrower and it is about whether the validation problem is
\emph{well posed}. For a classifier it is. The output space is finite and
enumerable, so a complete reference set can in principle be constructed. Errors are
label mismatches: atomic, countable, and unambiguous. The input space carries a
metric, so perturbation robustness can be stated as a radius and, for some model
classes, certified within it. Residual risk concentrates in distribution shift,
which is a monitorable quantity with established practice around it.

None of those four properties holds for an assertion. The output space is
$\mathcal{V}^{*}$, so there is no complete reference set to build. An error is not a
label mismatch: an assertion can be right in one jurisdiction and wrong in another,
adequate for a weaker claim than the one made, or true for reasons unconnected to
its evidence. The equivalence classes on the input are semantic rather than metric,
so there is no radius to certify and no coverage strategy against a class you cannot
enumerate. And the assertion makes a claim, so its correctness is relative to a
context that the input may never have specified.


\section{Hallucination as Ungrounded Assertion}
\label{sec:taxonomy}

The architecture of Section~\ref{sec:architecture} produces failures with a specific
shape: assertion is computed over surface form, with no evidence gate and no resolved
context, so output can be fluent and detached from anything at once. The received
vocabulary for those failures misdescribes them, and the misdescription has consequences
for what gets measured.

The term \emph{hallucination} is in universal use and is used in at least two
incompatible ways. In popular and much industry usage it is truth-conditional: a
hallucination is a false or fabricated output. In the natural-language-generation
literature it generally is not. There, hallucination means content unsupported by
the source, whether or not that content happens to be true, and the possibility of a
factually correct hallucination is explicitly recognised
\citep{maynez2020,ji2023}.

This paper adopts the second convention and claims no priority for it. Separating
support from truth is established work, and a definition that merely restated it
would contribute nothing. The definition below departs from \emph{both} conventions
somewhere else, and it is worth being exact about where, because that departure is
this paper's first contribution and it is narrower and more specific than
``hallucination is not falsity.''

The source-based definition takes the source as given. It asks whether the assertion
follows from the evidence supplied. It does not ask the prior question: whether that
evidence \emph{governs the decision at hand}. A clinical guideline may be authentic,
authoritative, correctly retrieved, accurately quoted, and genuinely entailing of the
assertion---and concern a different jurisdiction, a superseded period, or another
patient population. Every faithfulness metric scores that assertion as clean, because
the entailment holds. Definition~\ref{def:grounding} does not, because it requires
support to be \emph{applicable at a resolved context}.

That is the move: hallucination indexed to a decision context rather than to a
source. It is the first of this paper's two central claims. The second, developed in
Section~\ref{sec:canonical-form}, is that the same architecture makes accuracy
uncertifiable. They are connected, and Section~\ref{sec:unity} states how.

\subsection{Assertion}

I will talk about the deployed system $S$ rather than the model $M$ in isolation.
That is the thing being validated, and it is the thing whose failures have owners.

\begin{definition}[Assertion]
\label{def:assertion}
$\Out_S(p)$ holds iff system $S$, under input $x$, emits a claim with content $p$
at output probability at least $\theta$.
\end{definition}

The threshold $\theta$ separates commitment from mention: a system that lists $p$
among possibilities has not asserted it. It is defined over output probability,
measurable from the decoding distribution. Nothing here requires a view about what
$M$ believes, per Section~\ref{sec:premise}, and every predicate below is likewise
defined over instrumentable behaviour.

\subsection{Grounding}
\label{sec:grounding-def}

Grounding is the property that carries the paper, and it must mean considerably
more than that a document was retrieved or a citation was emitted.

\begin{definition}[Grounded evidence]
\label{def:grounded-evidence}
For an evidential set $E$, proposition $p$, and resolved context $c$:
\begin{equation}
  \begin{aligned}
  \Grounded(E, p, c) \;\leftrightarrow\;
    &\Authentic(E) \wedge \Versioned(E) \wedge \Authoritative(E, c)
      \wedge \Applicable(E, c)\\
    &\wedge\ \Supports(E, p, c) \wedge \Sufficient(E, p, c)
      \wedge \Traceable(E, p, c).
  \end{aligned}
  \label{eq:grounded}
\end{equation}
\end{definition}

\begin{definition}[Grounded assertion]
\label{def:grounding}
$\Grd_S(p \mid c)$ holds iff there is some $E \subseteq E_S$ in the system's
evidential base such that
\begin{equation}
  \Grounded(E, p, c) \ \wedge\ \Sensitive(S, p, E) \ \wedge\ \NoDefeater(E_S, p, c),
  \label{eq:grd}
\end{equation}
where $\Sensitive(S, p, E)$ holds iff, had $E$ not supported $p$, the assertion would
not have been produced, and
\begin{equation}
  \NoDefeater(E, E_S, p, c) \;\equiv\; \DefeaterSearch(E, E_S, p, c) = \emptyset .
  \label{eq:nodefeater}
\end{equation}
$\DefeaterSearch$ is a \emph{specified search} over the governed base, not a
quantification over its subsets. It returns evidence $E'$ that is applicable at $c$, of
standing at least equal to that of $E$, and that either supports $\neg p$ or undermines
$E$'s support for $p$:
\begin{multline}
  \DefeaterSearch(E, E_S, p, c) = \bigl\{\, E' \in \mathcal{Q}(E_S, c) \;:\;
    \standing(E', c) \succeq \standing(E, c)\\
    \wedge\ \bigl[\Supports(E', \neg p, c)
      \ \vee\ \text{undermines}(E', E, p, c)\bigr] \,\bigr\}.
  \label{eq:defeatersearch}
\end{multline}
\end{definition}

The seven conjuncts of \eqref{eq:grounded} are properties of the \emph{evidence};
$\Sensitive$ is a property of the \emph{system}; and $\NoDefeater$ is a property of
the \emph{evidential base as a whole}. All three kinds are required, and the
separation matters operationally because they are checked by different means and fail
for different reasons. Section~\ref{sec:grounding-context} takes the evidence
dimensions one at a time; here I note why the other two cannot be dropped.

\subsection{Why grounding needs a defeater condition}
\label{sec:why-defeater}

$\NoDefeater$ closes a hole that the existential quantifier in \eqref{eq:grd} would
otherwise open. Without it, grounding is satisfied whenever \emph{some} applicable
authoritative source supports $p$---even if another source of equal standing in the same
base supports $\neg p$. That is not a corner case. Two regulators may differ; a standard
and its national implementation may differ; a clinical guideline and a local protocol
may differ; a group policy and a subsidiary policy may differ. A system that retrieves
one side of such a disagreement and asserts it has not produced a grounded assertion. It
has produced an arbitrary selection from a contested record.

Three features of \eqref{eq:nodefeater} and \eqref{eq:defeatersearch} are deliberate,
and each answers an objection the unbounded form invites.

\paragraph{It is a search, not a quantifier.} $\mathcal{Q}(E_S, c)$ is the governed
retrieval over the base under the applicability predicates of $c$---the same retrieval
the system already performs, run in the contrary direction. Quantifying over subsets of
a corpus would be neither computable nor useful: in a large regulatory base something
can be read as supporting almost any negation. Bounding the condition to a specified
search makes it a control that a system can execute and a reviewer can re-execute, and
it aligns the definition with the architecture of
Section~\ref{sec:governed-rag}. The cost is honest and worth stating: a defeater outside
the search is a corpus-coverage failure rather than a loophole in the definition, and it
is detected the way coverage failures are always detected---by widening the search, not
by strengthening the logic.

\paragraph{Standing is an ordering, not a predicate.} $\standing(E, c)$ places evidence
in a partial order of authority that the institution defines: statute above guidance,
regulator above trade body, group policy above local practice except where local is
stricter. Contrary evidence of \emph{lower} standing is not a defeater---it is resolved
by precedence, which is what precedence is for. Only contrary evidence of equal or
incomparable standing defeats. This matters practically: the unordered version would
escalate every routine conflict, and a control that escalates routine cases is switched
off. It also locates a deliverable: the precedence order is a governed artefact owned by
the policy function, and its absence is why conflicts currently resolve by retrieval
order.

\paragraph{Defeat comes in two kinds.} The vocabulary is borrowed from work on
defeasible reasoning, where the distinction is standard. A \emph{rebutting} defeater
supports the contrary claim. An \emph{undercutting} defeater leaves the contrary claim
untouched and attacks the support relation instead: the instrument was withdrawn, its
scope was narrowed by a later provision, the passage was superseded in a revision the
retrieval did not reach, or the extraction misread a table. Undercutting defeat is
common in a governed corpus and invisible to a search for contradictions, because
nothing contradicts anything---the support simply no longer holds. Both belong in
\eqref{eq:defeatersearch}.

Where $\DefeaterSearch$ returns a result, the correct outcome is adjudication or
escalation rather than assertion. Section~\ref{sec:contested} treats this as a category
in its own right.

Every output of $M$ is caused by training, so causal provenance cannot be the
criterion---it would be satisfied by everything. And $\Grounded(E, p, c)$ alone is
satisfied whenever suitable evidence happens to exist somewhere in $E_S$, whether
or not the assertion had anything to do with it. A system that retrieves an
authoritative passage and then answers from parametric memory, coincidentally
agreeing with it, satisfies \eqref{eq:grounded} and is not grounded in any sense
an auditor should accept. $\Sensitive$ is what excludes this, and its operational
test is ablation: remove or contradict the putatively supporting evidence and
observe whether the assertion survives.

$\Sufficient$ closes a different hole, and one a data-quality reviewer will find
first. Evidence can be authentic, versioned, authoritative, applicable, and genuinely
supporting, and still address only part of what the claim requires. A rule imposing
five conditions is not satisfied by a source that speaks to three of them, and an
assertion resting on such a source is grounded in every dimension except adequacy.
Completeness is a standard data-quality dimension and its omission here would be
conspicuous.

The two conjuncts are independent rather than redundant, and the reason should be
stated because the redundancy objection is otherwise natural. $\Supports$ is a
\emph{positive} relation admitting degrees: the evidence favours $p$, at some strength,
as Section~\ref{sec:entails-supports} sets out. $\Sufficient$ is a \emph{coverage}
relation: the evidence addresses every condition the claim requires. Evidence can favour
a partial claim strongly---entail it, even---while covering three of five conditions, so
neither conjunct implies the other. Their remedies differ accordingly: inadequate support
is fixed by retrieving \emph{better}, insufficient coverage by retrieving \emph{more}.

\subsection{Five categories}

Let $A_i \in \{0,1\}$ record whether assertion $i$ is true under its resolved
context $c_i$, and $G_i \in \{0,1\}$ whether it is grounded in the sense of
Definition~\ref{def:grounding}. Crossing the two gives four categories.

\begin{table}[ht]
  \centering
  \begin{tabular}{@{}cc>{\raggedright\arraybackslash}p{2.4cm}>{\raggedright\arraybackslash}p{7.4cm}@{}}
    \toprule
    $A_i$ & $G_i$ & Category & Meaning \\
    \midrule
    1 & 1 & \textsc{Sound}
      & True assertion with authenticated, applicable, traceable evidentiary
        support \\
    \addlinespace
    1 & 0 & \textsc{Lucky}
      & True assertion lacking a sufficient auditable evidentiary basis \\
    \addlinespace
    0 & 1 & \textsc{Error}
      & Grounded but false. Decomposes by stage: the source is wrong; the passage
        was misread or mis-extracted; an inference step failed; a calculation
        failed; or the evidence was adequate for a weaker claim than the one made \\
    \addlinespace
    0 & 0 & \textsc{Fabrication}
      & False assertion without an auditable evidentiary basis \\
    \bottomrule
  \end{tabular}
  \caption{The four categories. Reference-based evaluation observes $A_i$ only.}
  \label{tab:taxonomy}
\end{table}

Because grounding requires applicability, a fifth category is analytically
distinct and, in regulated deployment, the most consequential.

\begin{table}[ht]
  \centering
  \begin{tabular}{@{}>{\raggedright\arraybackslash}p{5.6cm}>{\raggedright\arraybackslash}p{2.2cm}>{\raggedright\arraybackslash}p{5.4cm}@{}}
    \toprule
    Condition & Category & Meaning \\
    \midrule
    $\Authentic(E) \wedge \Authoritative(E, c') \wedge \neg\,\Applicable(E, c)$
      & \textsc{Misapplied}
      & Authentic and perhaps authoritative evidence that does not govern $c$:
        wrong jurisdiction, time period, population, policy, product, or decision
        use \\
    \bottomrule
  \end{tabular}
  \caption{The fifth category. Operationally this is a subset of
    $G_i = 0$, since Definition~\ref{def:grounding} requires applicability. It is
    kept separate because it identifies a different control failure with a
    different owner.}
  \label{tab:misapplied}
\end{table}

A sixth follows from $\NoDefeater$ for the same reason:

\begin{itemize}
  \item[\textsc{(vi)}] \textsc{Contested} --- the governed defeater search returns
        applicable evidence of equal or incomparable standing on the other side. No
        selection is grounded, whichever side the system asserts and whichever turns out
        to be right.
\end{itemize}

Like \textsc{Misapplied}, this is operationally a subset of $G_i = 0$: the no-defeater
condition of Definition~\ref{def:grounding} fails, so the assertion is ungrounded and
falls in the right-hand column of Table~\ref{tab:taxonomy}. It is kept separate for the
same reason---the control that fails is different, the remedy is different, and the
owner is different. There is no double counting in either case.

The relationship between Table~\ref{tab:misapplied} and Table~\ref{tab:taxonomy}
should be stated explicitly, since it is a place where a reader may reasonably
suspect double-counting. Grounding as defined requires applicability, so a
misapplied assertion is \emph{operationally ungrounded}: $G_i = 0$, and it falls in
the right-hand column. It remains analytically distinct because the underlying
control failure is different in kind. In cells (ii) and (iv) no evidence was
brought to bear at all. In the misapplied case authentic evidence was retrieved,
cited, and genuinely relied upon---and its scope did not match the assertion's
resolved context. The first is a generation failure. The second is a binding
failure. Work on the generation path will rarely find it, because the information that
would reveal the mismatch---which regime governs, which version was in force at
$t_v$---is metadata about the evidence rather than content of the assertion, and is
usually not in the context at all.

\subsection{The definition}
\label{sec:hall-def}

\begin{definition}[Hallucination, error, misapplied evidence]
\label{def:hallucination}
\begin{align}
  \Hall(p \mid c)       &\;=_{\mathrm{def}}\; \Out_S(p) \wedge \neg\,\Grd_S(p \mid c), \label{eq:hall}\\
  \Err(p \mid c)        &\;=_{\mathrm{def}}\; \Out_S(p) \wedge \Grd_S(p \mid c) \wedge \neg p, \label{eq:err}\\
  \Misapplied(p \mid c) &\;=_{\mathrm{def}}\; \Out_S(p) \wedge \Supports(E, p, c')
                           \wedge \Authentic(E) \wedge \neg\,\Applicable(E, c),
                           \label{eq:misapplied}\\
  \Contested(p \mid c)  &\;=_{\mathrm{def}}\; \Out_S(p) \wedge \neg\,\NoDefeater(E, E_S, p, c).
                           \label{eq:contested}
\end{align}
\end{definition}

Note what \eqref{eq:hall} omits: \textbf{any condition on the truth of $p$.} A
hallucination is an assertion unsupported by applicable evidence, and that is all.
Truth and falsity distinguish successes from errors; they do not distinguish
hallucination from sound assertion.

Two consequences follow immediately, and both invert the standard usage. A false
assertion faithfully reporting an applicable source is \emph{not} a hallucination;
it is an error, and the fault lies in the evidence. A true assertion resting on
nothing \emph{is} a hallucination, and remains one however often it turns out
right. The next three subsections defend each of these, and the defence is in each
case an argument about remedies rather than about words.

The definition also has a property that neither convention has: it is
\emph{evaluable at the point of assertion}, without a reference answer. That is the
subject of Section~\ref{sec:auditable}, and it is why this is an engineering claim
rather than a terminological preference. A truth-conditional definition can only be
applied offline, against known answers. A source-based definition can be applied at
runtime but only relative to whatever was retrieved. Definition~\ref{def:grounding}
is applied at runtime against the evidence \emph{and} the context, which is what a
gate needs.

\subsection{What is claimed as new, and what is not}
\label{sec:novelty}

Groundedness, provenance, factuality, attribution, and contextual relevance all have
substantial literatures, and nothing here is offered as a discovery in any of them.
The claim is narrower, and stating it precisely is worth a paragraph because a claim
pitched wider would be false.

This paper does not treat hallucination as synonymous with factual error, and does
not claim priority for that separation. It formalises the relevant failure as the
absence of an auditable relation between an assertion, an authenticated and
\emph{contextually applicable} evidentiary basis, and a verification pathway that an
institution has authorised. Four components of that formalisation are, so far as I am
aware, not present in existing definitions:

\begin{enumerate}
  \item \textbf{Applicability as a conjunct.} Support is indexed to
        $c = \{j, t, f, u\}$, so evidence that entails the assertion but does not
        govern the case fails the definition. Faithfulness metrics take the source as
        given and cannot express this condition.
  \item \textbf{Dependence rather than entailment.} $\Sensitive$ requires the
        assertion to have been produced \emph{because of} the evidence, not merely to
        be consistent with it. A system that retrieves an authoritative passage and
        answers from parametric memory satisfies entailment and fails this.
  \item \textbf{Authority and version as conditions on support.} Whether a source has
        standing for this decision class, and which version of it applied at $t$, are
        treated as preconditions of valid support rather than as metadata.
  \item \textbf{Defeat by contrary authority.} $\NoDefeater$ makes grounding a
        property of the evidential base rather than of a selected subset of it, so an
        assertion drawn from one side of a live disagreement between applicable
        authorities is not grounded. Faithfulness and attribution metrics evaluate an
        assertion against the source it cites and have no way to express this.
  \item \textbf{The independence result.} Proposition~\ref{prop:independence}
        establishes that no accuracy figure constrains the resulting rate, which is
        what makes the redefinition consequential for validation rather than
        expository.
\end{enumerate}

Taken together these separate four cases that pooled accuracy and conventional
source-faithfulness metrics conflate: lucky truth, grounded error, fabrication, and
contextual misapplication. The compact form of the claim is that factual accuracy is
insufficient for institutional reliance, because an assertion must also possess a
reconstructable and contextually applicable evidentiary lineage---and that a metric
which cannot distinguish those four cases cannot support a reliance decision.

% TODO(priority): I have not been able to establish exhaustively that no existing

\subsection{Cell (iii): a faithful report of a bad source is an error, not a fabrication}
\label{sec:cell-iii}

A corpus contains, consistently and across several reference works, an incorrect
effective date for a regulatory provision. The system reports that date. Ablation
confirms the assertion is driven by those documents: remove them and it stops.

Reference-based evaluation scores this as a hallucination, because the output is
false.
That classification destroys information the organisation needs. The system did
what a reporting component is supposed to do: it reported what its evidence said,
through a mechanism sensitive to that evidence. The defect is in the evidence.

The distinction is not terminological, it is budgetary. Cell (iii) is addressed by
work on the evidential base---source selection, authority ranking, curation,
retrieval from systems of record rather than secondary literature. Cell (iv) is
addressed by work on the generation path---evidence gating, abstention policies,
refusal thresholds. Different teams, different skills, different funding cycles, as
Section~\ref{sec:three-teams} develops. A metric that pools them reports a number
on which no one can act.

There is a normative point as well. We do not say that a careful analyst
fabricated when she accurately reported a source that later proved wrong. We say
she was misled, and we look at her sources. Machine assertion should be held to the
same distinction, in both directions.

\subsection{Cell (ii): a lucky right answer is not evidence about the system}
\label{sec:cell-ii}

Now the reverse. The system is asked for a figure absent from its evidential
base---the headcount of a small subsidiary. The generation path, running on surface
regularities about how entities of this type and sentences of this genre continue,
produces $8{,}200$. That happens to be right.

Reference-based evaluation scores this as a success. Every property that makes
fabrication dangerous is nonetheless present, and the precise defect is worth
stating because it is what makes the case more than a curiosity: had the true
figure been $3{,}400$, the same process would have produced the same output. The
assertion does not depend on the fact it reports. It is not a measurement; it is a
draw from a distribution that happened to land on the answer.

The consequence for validation is the one that matters. A user who receives a
cell-(ii) output, checks it, and finds it correct has acquired \emph{no} evidence
about the system. The check confirmed the output, not the process, and the process
will behave identically on the next question, where the draw may land elsewhere.
Lucky answers are not repeatable performances. Grounding is---and that asymmetry
is the argument of Section~\ref{sec:auditable}, arriving early.

This also explains why a benchmark that rewards cell (ii) is actively harmful
rather than merely uninformative. Scoring lucky answers as successes applies
optimisation pressure toward confident guessing and away from abstention, which by
Section~\ref{sec:rlhf} is a direction the training regime already favours. The
metric compounds a bias the architecture supplies.

\subsection{Cell (v): authentic is not applicable}
\label{sec:cell-v}

The category pooled metrics hide, and in my experience the most common serious
failure in regulated deployment.

The system asserts that a treatment is indicated for a condition. The assertion is
driven by clinical guidance that says exactly that, and ablation confirms it. But
the guidance concerns an adult population and the case is paediatric. Or a system
asserts a capital treatment on the basis of a rule that governs a different
jurisdiction. Or it cites a policy that was superseded before the decision date. Or
it applies a product rule to a different product class.

Without the context index these are cell (iii): grounded, false, therefore charged
to a defective evidential base. That verdict is wrong twice over. The base was not
defective---the guidance is accurate for its population, the foreign rule is
correctly stated, the superseded policy was valid when issued. And the implied
remedy is wrong: curating the corpus will not help, because nothing in it is
incorrect.

What failed is applicability, and the remedy is context resolution: establishing
$c$ before answering, checking currency against $t$, checking jurisdiction,
material facts, and intended use, and abstaining or qualifying when the applicable
evidence has not been identified. These examples are illustrations of the failure
shape rather than reports of measured incidents, and
Section~\ref{sec:grounding-context} develops the underlying generalisation error.
\subsection{Cell (vi): a contested record is not a grounded one}
\label{sec:contested}

The failure that the existential quantifier hides, and the one most likely to be
misclassified as an error.

Two authoritative sources in the base disagree. A group accounting policy and a
statutory requirement in one jurisdiction reach different treatments. A clinical
guideline and a local formulary differ on a first-line therapy. A standard and its
national implementation diverge on a threshold. The system retrieves one, and asserts
it. Ablation confirms the assertion rests on that source; the source is authentic,
versioned, authoritative, applicable, sufficient, and traceable.

Every dimension of \eqref{eq:grounded} passes, and the assertion is not fit for
reliance. What the system has done is resolve a policy conflict by retrieval order,
which is to say arbitrarily, and it has produced a record showing a single clean
citation with no indication that the question was contested at all. That is worse than
an unsupported answer, because the audit trail actively conceals the problem: a
reviewer sees one authoritative source and an assertion that follows from it.

Note that the defect is independent of which side is correct. If the retrieved source
happens to state the treatment a reviewer would have chosen, nothing has been
established---the selection was still arbitrary, and the next case may go the other
way. This is the lucky cell again, at the level of source selection rather than token
generation.

The remedy is neither corpus work nor generation work. It is \emph{adjudication}: a
rule that determines which authority governs when authorities conflict, and an
escalation path for cases the rule does not settle. That is a policy artefact, and its
absence is why contested cases surface as apparently well-grounded assertions.

The categories exist to route work, and this is what they are for. Four remedies,
across three owners:
\label{sec:remedies}

\begin{table}[ht]
  \centering
  \begin{tabular}{@{}>{\raggedright\arraybackslash}p{2.6cm}>{\raggedright\arraybackslash}p{4.0cm}>{\raggedright\arraybackslash}p{4.4cm}>{\raggedright\arraybackslash}p{2.3cm}@{}}
    \toprule
    Category & Diagnosis & Remedy & Typical owner \\
    \midrule
    \textsc{Error}
      & applicable evidence, but wrong
      & source authority, curation, systems of record
      & data \\
    \addlinespace
    \textsc{Misapplied}
      & evidence inapplicable at $c$
      & context resolution, currency and scope binding
      & policy / risk \\
    \addlinespace
    \textsc{Contested}
      & applicable authority on both sides
      & precedence rules, adjudication, escalation path
      & policy / risk \\
    \addlinespace
    \textsc{Lucky}, \textsc{Fabrication}
      & no evidential connection
      & evidence gating, abstention, refusal thresholds
      & platform \\
    \bottomrule
  \end{tabular}
  \caption{A single pooled hallucination rate cannot distinguish these rows, and
    therefore cannot tell an organisation which of four budgets to spend.
    Section~\ref{sec:three-teams} argues that the division of ownership is itself
    a principal source of failure.}
  \label{tab:remedies}
\end{table}


\section{Accuracy and Grounding Are Independent}
\label{sec:independence}

The redefinition of Section~\ref{sec:hall-def} is not merely a useful
reclassification. It supports a formal result about what aggregate accuracy can and
cannot tell an assessor---and the result is what makes the redefinition consequential
rather than semantic.

\subsection{The result}

For a population of $n$ assertions, define
\begin{equation}
  \mathrm{Accuracy} = \frac{1}{n}\sum_{i=1}^{n} A_i,
  \qquad
  \mathrm{GroundingRate} = \frac{1}{n}\sum_{i=1}^{n} G_i.
  \label{eq:rates}
\end{equation}

\begin{proposition}
\label{prop:independence}
Let $\alpha, \gamma \in [0,1]$ with $n\alpha$ and $n\gamma$ integral. There is a
population of $n$ assertions with $\mathrm{Accuracy} = \alpha$ and
$\mathrm{GroundingRate} = \gamma$. Aggregate accuracy therefore neither identifies
nor bounds the grounding rate, and the grounding rate neither identifies nor bounds
aggregate accuracy.
\end{proposition}

\begin{proof}
Let $n_{\mathrm{sound}}, n_{\mathrm{lucky}}, n_{\mathrm{err}}, n_{\mathrm{fab}}$
be the counts in the four cells of Table~\ref{tab:taxonomy}, so that
$n\alpha = n_{\mathrm{sound}} + n_{\mathrm{lucky}}$ and
$n\gamma = n_{\mathrm{sound}} + n_{\mathrm{err}}$. Choose any integer
\[
  n_{\mathrm{sound}} \in
  \bigl[\max(0,\, n\alpha + n\gamma - n),\ \min(n\alpha,\, n\gamma)\bigr],
\]
an interval containing an integer for every admissible $\alpha, \gamma$, and set
$n_{\mathrm{err}} = n\gamma - n_{\mathrm{sound}}$,
$n_{\mathrm{lucky}} = n\alpha - n_{\mathrm{sound}}$, and
$n_{\mathrm{fab}} = n - n\alpha - n\gamma + n_{\mathrm{sound}}$. Each is a
non-negative integer by the choice of interval, and the four sum to $n$.
\end{proof}

The construction requires only that grounding not entail truth, which cell (iii)
supplies: an assertion can rest on authentic, applicable, authoritative evidence
and still be false, through misreading, invalid inference, or evidence adequate for
a weaker claim than the one asserted.

The result is arithmetically simple and strategically significant, and the two
should not be confused. For a population of generated assertions, any feasible
aggregate accuracy can coexist with any feasible grounding rate, because the
population may contain varying proportions of grounded-true, ungrounded-true,
grounded-false, and ungrounded-false assertions. Consequently pooled accuracy
cannot identify or bound the frequency of auditable grounding.

One limitation should be stated before the corollary, because a reader is entitled to
ask for it. The construction establishes \emph{logical} independence: no accuracy
figure entails anything about the grounding rate. It does not establish that the two
are empirically uncorrelated in deployed systems, and they may well be---a system
built to ground its assertions will plausibly be more accurate too. The claim is about
what an accuracy measurement \emph{licenses}, not about what is typical. That is
sufficient for the argument, because a validation regime rests on entailment rather
than on correlation: an assessor who accepts an accuracy figure as evidence of
grounding is relying on a relation that does not hold, whatever the empirical
association happens to be in the population they did not measure.

\begin{corollary}
\label{cor:no-bound}
A system measured at $95\%$ accuracy is consistent with a grounding rate of $0$,
and a system measured at $40\%$ accuracy is consistent with a grounding rate of
$1$.
\end{corollary}

\subsection{What it implies for a validation programme}

What follows is a claim about validation programmes rather than about models. A
programme that reports accuracy without measuring per-assertion grounding cannot
distinguish between radically different reliance profiles---between a system whose
correct answers rest on identifiable authoritative evidence and one whose correct
answers are draws that happened to land well. Those two systems have the same
score and different futures, and the difference appears the moment the question
distribution moves.

Two further readings matter for anyone choosing between releases. Since the two
rates move independently, an intervention that raises accuracy tells you nothing
about grounding; and optimising against reference answers applies pressure toward
confident guessing, which raises accuracy by moving mass into the lucky cell where
$G_i = 0$. A leaderboard improvement is fully consistent with a system that has
become more willing to assert without support.


\section{Grounding, Context, and Applicability}
\label{sec:grounding-context}

Definition~\ref{def:grounded-evidence} listed seven conditions on evidence without
explaining them. This section takes them in turn and identifies the generalisation error
that makes applicability the hardest of the seven to enforce. The resolved context $c$
that four of them depend on was defined in Section~\ref{sec:resolved-context}.

\subsection{Illicit universalization}
\label{sec:illicit}

Relation \eqref{eq:no-transfer} names a failure that deserves its own label,
because treating it as ordinary inaccuracy leads to the wrong remedy.

\emph{Illicit universalization} is the treatment of a claim that holds under one
jurisdiction, time, population, policy, or use as though it held generally, because
the relevant context was never resolved. The claim is not fabricated. It is not
even, in its own frame, wrong. It has been detached from the frame that made it
assessable, and then asserted into a frame where its truth value is simply a
different question.

This is worth insisting on because it reframes a large fraction of what gets
reported as hallucination. Many apparent failures are not invented facts. They are
unresolved-reference failures, temporal-validity failures, jurisdiction failures,
or applicability failures: information that is correct somewhere, stated without a
valid binding to the context in which reliance is proposed. The diagnostic question
is not ``is this true?'' but ``true \emph{where}, \emph{when}, \emph{for whom}, and
\emph{for what use}?''---and a system that cannot represent the question cannot be
said to have answered it.

The data-architecture analogue is exact, and worth stating because it makes the
failure recognisable rather than novel. A measure without its dimension keys is not
a fact; it is a number. Revenue is meaningless without period, entity, and currency,
and no data architect would accept it into a warehouse unqualified. Illicit
universalization is the same defect at the level of an assertion: a claim detached
from the dimensions that make it assessable, presented as though those dimensions
were universal rather than merely unstated. What ordinary dimensional modelling
enforces by schema, these systems do not enforce at all, because by
\eqref{eq:signature} there is no schema to enforce it with.

There is a structural reason this failure is native to the architecture rather than
incidental. By \eqref{eq:signature} the system's input is a token sequence, and a
token sequence carries no context index. A question that omits its jurisdiction is
not, to the system, an incomplete question; it is a complete string. Nothing in
\eqref{eq:objective} rewards noticing that a binding is missing, and by
Section~\ref{sec:rlhf} preference optimisation actively rewards answering anyway.
The system is not failing to resolve context. It has no representation of context
to resolve, and no mechanism that could register the omission.

\subsection{The seven dimensions}
\label{sec:six-dimensions}

\paragraph{Authentic.} The source is what it purports to be. In practice this means
retrieval from a controlled corpus whose contents have known origin, not from a
scrape whose provenance is a URL. Authenticity is a property an organisation
establishes by ingestion controls, and it is the dimension most familiar from
ordinary data governance.

\paragraph{Versioned.} The source version, its effective date, and the corpus state
$t_k$ at retrieval can all be identified. A regulation without an effective date is
not evidence about what was required at a decision date; it is evidence about what
some document said at some point. Both timestamps are needed, for different purposes:
the effective date establishes validity at $t_v$, and $t_k$ establishes what the
system could have seen. This is the dimension that permits replay, and without it the record itself cannot
support reconstruction: whatever an investigator eventually pieces together from
archives and change logs, it will not be a reconstruction from the system's own
evidence.

\paragraph{Authoritative.} The source has the standing required for the decision
type. A vendor summary of a regulation and the regulation itself may agree on
content and differ entirely in standing, and which one is required is a function of
the decision, not of the content. Authority is a relation between a source and a
decision class, not a scalar attached to a document---which is why an authority
ranking that is global rather than use-conditional will misclassify.

\paragraph{Applicable.} The source governs or materially bears on $c$. Read
component-wise, this requires the source to govern in jurisdiction $j$, to be in force
at valid time $t_v$---interval containment, per
Section~\ref{sec:resolved-context}, so neither not-yet-enacted nor
already-superseded---to bear on the material facts $f$, and to cover the intended use
$u$. This is where \eqref{eq:no-transfer} bites, and it is the dimension standard
retrieval most reliably fails, for reasons developed in
Section~\ref{sec:retrieval-insufficient}.

\paragraph{Supports.} The source establishes $p$ at the strength asserted.
Section~\ref{sec:entails-supports} distinguishes the grades of support this covers.

\paragraph{Sufficient.} The source addresses every condition the claim requires,
not merely some of them. A rule imposing five conditions is not satisfied by evidence
speaking to three, and an assertion resting on partial evidence is grounded in every
other dimension while being inadequate in the one that matters. Sufficiency and
support are separated because their remedies differ: insufficient support is fixed by
retrieving \emph{more}, inadequate support by retrieving \emph{better}.

\paragraph{Traceable.} A reviewer can reconstruct the evidence-to-assertion
relation: which passage, which version, retrieved at which corpus state, supporting
which claim. Traceability is a property of the record rather than of the evidence, and
it is what turns the other six into an auditable finding rather than an assertion
about the system's internals.

Note that four of the seven---$\Authoritative$, $\Applicable$, $\Supports$, and
$\Sufficient$---are relations to $c$ rather than properties of the document. A corpus
can be perfectly authentic, fully versioned, and completely traceable while failing
all four, and a governance programme that treats grounding as a corpus property will
build exactly that system. The no-defeater condition of \eqref{eq:nodefeater} is
stronger still: it is a relation to the \emph{entire base}, so no examination of a
single retrieved source can establish it.

\subsection{Entailment and support}
\label{sec:entails-supports}

$\Supports(E, p, c)$ covers a range, and collapsing it is a common source of
overclaiming. At the strong end,
\begin{equation}
  \Entails(E, p, c) \;\Rightarrow\; \Supports(E, p, c),
  \label{eq:entails}
\end{equation}
where $E$ together with the governing rules of $c$ settles $p$. A capital
calculation that follows deterministically from a rule and a set of positions is of
this kind, and so is a policy question whose answer is stated in the policy.

Most enterprise questions are weaker. Evidence may make $p$ substantially more
likely without settling it; it may support a qualified version of $p$ but not the
unqualified assertion; it may support $p$ for a class of cases of which $c$ is
plausibly but not demonstrably a member. None of this is a defect. It is the normal
condition of decision-making under real evidence.

The requirement is not that every assertion be entailed. It is that the kind and
strength of support be \emph{explicit and proportional to the consequence of
reliance}. A system that reports ``supported'' without distinguishing entailment
from suggestive correlation has not given a reviewer what they need, and has
created the specific failure in which an assertion adequate for one decision class
is relied upon for another. This is cell (iii) in its most defensible form: the
evidence was authentic, applicable, and authoritative, and it was adequate for a
weaker claim than the one asserted.

\subsection{Why retrieval and citation are not enough}
\label{sec:retrieval-insufficient}

The anticipated objection to this entire section is that retrieval-augmented
generation \citep{lewis2020} already addresses it. Retrieval does help, materially,
and no argument here counts against using it. It does not establish grounding in
the sense of Definition~\ref{def:grounded-evidence}, and the gap is specific.

Standard retrieval ranks passages by semantic similarity to the query and places
the top results in the context window. Consider what that procedure establishes
against the six dimensions. It can establish authenticity, if the corpus is
controlled. It can establish traceability, if the retrieval event is logged. It
establishes nothing about versioning unless version metadata is carried and used,
nothing about authority unless authority is represented and conditioned on the
decision type, and---critically---nothing about applicability, because semantic
similarity to the query is not applicability to the context.

That last point deserves to be stated as a slogan, because it is the practical core
of this paper:

\begin{quote}
Retrieval is not applicability. Citation is not entailment. And authoritative
evidence is not automatically governing evidence for the case at hand.
\end{quote}

The nearest existing measurement is worth naming, since retrieval evaluation does
score something adjacent. Context-relevance metrics assess whether retrieved evidence
is relevant \emph{to the query}. Applicability asks whether it \emph{governs the
case}. A paediatric question that retrieves adult guidance scores as highly
relevant---the topic matches exactly---and the guidance does not govern. Relevance to a
query and authority over a decision are different relations, and only the first is
measured.

Each clause of the slogan names a distinct substitution that systems make routinely. A
passage
about paediatric dosing and a passage about adult dosing are near-neighbours in
embedding space; that is what makes similarity search work, and it is exactly why
similarity search will surface the wrong one. A citation demonstrates that a
document was consulted, not that it supports the claim attached to it, and
citation-to-claim alignment is a separate check that standard pipelines do not
perform. And a source may be the most authoritative document in the corpus and
still not govern this decision.

Nor does $\Sensitive$ come for free. Placing evidence in the context raises the
probability of consistent continuations; it does not make the assertion conditional
on the evidence having supported it. By Section~\ref{sec:missing} retrieved text is
more tokens, not a precondition, and a system may retrieve an authoritative passage
and then answer from parametric memory, agreeing with it or not.

The conclusion is not that retrieval fails. It is that retrieval supplies one of
six dimensions reliably, two more if metadata is carried and logged, and none of
the three that are relations to $c$. The remaining work is architecture, and
Section~\ref{sec:governed-rag} sets out what it consists of.


\section{There Is No Guaranteed Canonical Form}
\label{sec:canonical-form}

One qualification belongs in the first paragraph, because the heading is stronger than
what follows and the difference is the point. Training on paraphrase-rich text does
apply pressure toward treating equivalent formulations alike, and embeddings do place
many paraphrases near one another. What is absent is not similarity; it is a
\emph{guarantee}---an enforced, inspectable, versionable canonical form of the kind
every other component in the stack supplies. A learned tendency and an enforced
invariant differ in exactly the way that matters for validation, and the rest of this
section is about that difference.

One property does most of the damage, and it is worth naming before it is formalised.
Call it \emph{form invariance}: materially equivalent requests should produce equivalent
decision-relevant outcomes. It is the property every other component in an enterprise
stack supplies without being asked, it is the property this architecture does not supply,
and its absence is what breaks validation. Section~\ref{sec:invariants} states it as
invariant I1 among the others; this section explains why it fails, what follows for a
validation programme, and why the obvious remedy does not work. The failure has a
one-sentence statement: the system is keyed on the surface form of its input, not on the
content.

\subsection{The map and what it does not respect}

Section~\ref{sec:what-computed} established the signature: $S$ maps
$\mathcal{V}^{*}$ to distributions over $\mathcal{V}^{*}$. The domain is token
sequences.

Now consider the equivalence relation that validation cares about. Write
$x \sim_c x'$ when $x$ and $x'$ encode \emph{materially equivalent} requests under
the same resolved context $c$---same question, same governing frame, same material
facts, same intended use. This relation partitions the input space. What a test
suite needs is that decision-relevant behaviour be constant on those classes:
\begin{equation}
  x \sim_c x' \;\Longrightarrow\; \text{EquivalentDecisionRelevantOutcome}(x, x', c).
  \label{eq:form-invariance}
\end{equation}

Two things about \eqref{eq:form-invariance} are deliberate. It is stated over
\emph{material} equivalence under a resolved context, not over semantic
equivalence in the abstract, because the equivalence that matters is domain- and
context-relative. And it requires equivalent decision-relevant \emph{outcomes},
not identical text. Nobody needs the system to produce the same sentence twice.
What an institution needs is that materially equivalent requests not yield
incompatible assertions, incompatible evidence bases, or incompatible
decision-relevant conclusions without an identifiable reason.

Nothing in the architecture delivers \eqref{eq:form-invariance}. There is no
normalisation pass, no canonical form, no quotient map. By
Section~\ref{sec:what-computed} the tokeniser is lexical, the embedding is of
tokens, and no later stage brings two encodings of one claim to a common form. Two
members of a $\sim_c$ class are two different points in $\mathbb{R}^{n \times d}$
with different distributional histories, and \eqref{eq:objective} contains no term
penalising divergence between them. The map is free to send them anywhere, and
sometimes does.

\subsection{Every other component in the stack does better}

The failure is easy to overlook because it is unlike anything else in an enterprise
architecture.

\begin{itemize}
  \item A database normalises before it compares. Results do not depend on whether
        a predicate was written \texttt{NOT (a OR b)} or \texttt{(NOT a) AND (NOT
        b)}; the planner canonicalises, and equivalence is decidable in the
        algebra. The comparison should not be overdrawn: equivalence for the full
        query language is not decidable either, and floating-point aggregation can
        differ across parallel plans for the same reasons discussed in
        Section~\ref{sec:determinism}. What a database supplies is a canonical form
        for the fragment that matters and a documented statement of where determinism
        is and is not guaranteed. Neither is available here.
  \item A compiler canonicalises before it optimises. \texttt{x*2} and
        \texttt{x+x} converge.
  \item A cache keyed correctly hashes normalised content, not raw bytes. Keying a
        cache on the unnormalised string is a familiar bug, and its symptom is
        familiar too: identical requests, different answers, no error anywhere.
\end{itemize}

The last case is the closest, and it is barely an analogy. The system behaves like a
cache keyed on the raw string: materially identical requests are distinct keys, and
nothing in the pipeline would notice.

The deeper version of the point is architectural, and it is the one that should
concern anyone who has built a system of record. A canonical record requires stable
identity, a normalised or governed representation, declared dependencies, lineage,
and reproducible derivation. A model's latent representation has none of these
properties: it has no stable identity for an entity across paraphrases, no
inspectable normalised form, no declared dependency on a source, and no reproducible
derivation path from evidence to output. It is therefore not a candidate system of
record, whatever its accuracy. Few architects would assert that it is; the error is
subtler and commoner, and consists in treating model \emph{output} as authoritative
without requiring the lineage that any other authoritative source in the enterprise
would have to supply.

Insofar as materially equivalent formulations \emph{are} treated alike---and often
they are---that is a contingent achievement of training coverage. The training
distribution contained many paraphrase pairs, so the map is approximately constant
on well-covered classes. This degrades exactly where prediction says it should: as
formulations become rarer, longer, more domain-specific, more like the queries a
specialist actually asks. The behaviour is best where the guarantee is needed
least.

\subsection{What the measurements show}

The prediction is testable and has been tested.

The cleanest case is format sensitivity. Holding semantic content fixed
and varying only formatting choices---separators, casing, spacing, the shape of an
enumeration---produces large spreads in task performance \citep{sclar2024}. This is
the result to show a sceptical colleague: the content of the request is unchanged by
construction, the variations are choices a prompt author makes arbitrarily and does
not record, and behaviour moves anyway. Aggregate performance cannot shift unless
individual answers changed, so a spread across formats is a lower bound on the rate
at which materially equivalent requests receive different treatment.

Framing sensitivity gives the direction, and is the second clean case. When a prompt signals what the user
would prefer to hear, assertion shifts toward it \citep{sharma2024}. The content
of the question is unchanged; the assertion is not.

Consistency under paraphrase, measured directly, falls well below what a
content-directed process would produce \citep{elazar2021}, and models will
contradict themselves within a single generation \citep{mundler2024}. The second
of these bears on I7 as well as I1: the instability is not only across sessions
and phrasings but inside one continuous output, where any notion of a maintained
position would have to apply if it applied anywhere.

\subsection{Why ``train a normaliser'' is not the answer}
\label{sec:normaliser}

The natural engineering response is to canonicalise the input: put a normaliser in
front of the model and the problem goes away. It does not, for five reasons that
compound.

\paragraph{Material equivalence is context-sensitive.} Whether two requests are
equivalent depends on $c$. ``Is this reportable?'' asked of two transactions
identical but for jurisdiction are not materially equivalent requests, and no
context-free normaliser can know that.

\paragraph{It depends on domain ontology.} Deciding that two phrasings pick out the
same product class, the same patient population, or the same instrument type
requires the domain's own identity conditions. That is ontology work, not string
work.

\paragraph{Paraphrases carry material distinctions.} In regulated language, surface
differences are frequently the point. ``Indicated for'' and ``may be used in'' are
not paraphrases. ``Shall'' and ``should'' are not paraphrases. A normaliser
aggressive enough to collapse genuine paraphrases will collapse these too, and the
failure will be silent and consequential.

\paragraph{General semantic equivalence is not reliably decidable.} Unlike
relational algebra, natural language offers no decision procedure for equivalence,
and any learned approximation inherits the properties of the thing it approximates,
including sensitivity to form.

\paragraph{A latent representation is not an inspectable canonical record.} Even a
normaliser that worked would produce an internal representation that no reviewer
can read, version, or diff---which fails the requirement it was introduced to
satisfy.

The conclusion is not that normalisation is useless. Structured intake, controlled
vocabularies, and typed request objects all help, substantially, and belong in the
architecture. The conclusion is that the remedy is not \emph{semantic}
normalisation of free text. It is explicit context binding, structured domain
representations where they can be built, governed source corpora, deterministic or
inspectable policy controls where the decision permits them, and a preserved
evidence-and-policy path for every consequential assertion. That is
Section~\ref{sec:governed-rag}.

\subsection{The validation consequence}
\label{sec:validation-consequence}

Now the argument this paper exists to make. State the validation inference
explicitly and it fails visibly.

Let $T = \{(x_1, a_1), \ldots, (x_n, a_n)\}$ be a test suite: inputs paired with
acceptable outputs. Suppose $S$ passes---for each $i$, $S(x_i)$ is acceptable. The
inference validation needs is:

\begin{quote}
For every $x$ that a user might actually send, if $x \sim x_i$ for some tested
$x_i$, then $S(x)$ is acceptable.
\end{quote}

That inference is valid only if \eqref{eq:signature} is constant on $\sim$-classes,
which is invariant I1. I1 fails. So the inference is invalid, and the test suite
does not license the conclusion drawn from it.

Three aggravating features, each of which an engineer will recognise as worse than
the bare failure.

\paragraph{The classes are unbounded.} You cannot enumerate the paraphrases of a
question, so you cannot fix the problem by testing more members of each class.
Coverage is not a strategy against an infinite equivalence class; the best you can
do is sample it, which returns you to sampling.

\paragraph{The failures are silent.} An I1 violation raises nothing. There is no
exception, no warning, no anomaly in the logs. The system returns a fluent,
well-formed, differently-wrong answer with the same confidence as the right one.
Contrast a schema violation or a type error, which announce themselves.

\paragraph{Repeatability fails too.} I7's failure means you cannot even rely on
the tested string. The same input can produce a different assertion later in the
same session, as context accumulates and shifts the conditioning. So the guarantee
does not merely fail to extend to paraphrases; it does not securely cover the
tested case either.

\subsection{Why this makes accuracy the wrong thing to certify}

Put the two halves together.

Accuracy is estimated by sampling: draw inputs, compare outputs to expected
answers, report the proportion correct. The estimate is a statement about the
distribution the inputs were drawn from. Using it to certify deployment requires
that the estimate transfer to the inputs that will actually arrive---and I1's
failure is precisely the claim that it does not transfer across semantic
equivalence, which is the dimension along which real inputs differ from test
inputs.

So accuracy is a property of a sample, and the sample does not generalise in the
direction validation needs. This is not a claim that the number is wrong. The
number is a correct measurement of something; it is a correct measurement of how
the system performed on those strings.

Notice, finally, what this does to the cell-(ii) case of
Section~\ref{sec:cell-ii}. A lucky assertion is one that would have been produced
whatever the fact had been. Verifying it tells you about the draw, not the
process. I1's failure generalises that observation from lucky assertions to
\emph{all} assertions: any verified output tells you about the draw. Cell (ii) was
the special case where the disconnection is total; non-repeatability is the same
disconnection showing up as instability across the input space.

That is the link between the two halves of this paper. The taxonomy of
Section~\ref{sec:taxonomy} says that a correct answer can be produced for no
reason. This section says that even a correct answer produced for a good reason
carries no guarantee of recurrence, because nothing holds the system to the
content rather than the string. Both point at the same missing property, and the
next section names it.

\subsection{Why the two claims are one}
\label{sec:unity}

Section~\ref{sec:taxonomy} redefined hallucination as ungrounded assertion. This
section argued that there is no canonical form, so a passing test establishes
nothing about a materially equivalent input. The two are the same architectural
fact seen from two sides, and it is worth saying how, because each makes the other
consequential.

Why do these systems hallucinate in the redefined sense? Because assertion
probability is computed over formulations. A string can be made probable by surface
statistics with no route through evidential support for its content, so assertion
and evidence come apart---which is exactly $\neg\,\Grd_S(p \mid c)$ while
$\Out_S(p)$ holds. Ungrounded assertion is what the absence of a content-level
representation looks like from the output side.

And why does the absence of a canonical form matter, rather than being an
interesting curiosity about paraphrase robustness? Because of what
Section~\ref{sec:cell-ii} showed about a lucky assertion: verifying it tells you
about the draw, not the process. Form-invariance failure generalises that from
lucky assertions to \emph{all} assertions. Any verified output tells you about the
draw, because nothing holds the system to the content rather than the string. The
lucky cell was the special case where the disconnection is total; form-invariance
failure is the same disconnection showing up as instability across the input space.

So the redefinition identifies the defect and the canonical-form argument explains
why you cannot test your way out of it. Together they force the conclusion of
Section~\ref{sec:auditable}: the property to measure is not whether assertions turn
out right, but whether they were grounded when made.


\section{System Invariants}
\label{sec:invariants}

Sections~\ref{sec:architecture} and~\ref{sec:canonical-form} argue from
architecture. This section argues from contract, which is the form a reviewer can
check without agreeing with me about mechanism.

If you were specifying a component whose assertions you intended to rely on, you
would write down invariants. Two groups are needed, and the distinction between
them is the paper's practical claim in compressed form. The first group states
properties of assertion \emph{behaviour}: things you would assume without thinking
about them from a database, a rule engine, or a calculation service. All of them
fail, and none can be made to hold by work on $M$. The second states properties the
\emph{system} must supply. None of them holds by default, and all of them can be
engineered.

\subsection{Behavioural invariants: what does not hold}

\begin{table}[ht]
  \centering
  \small
  \begin{tabular}{@{}>{\raggedright\arraybackslash}p{0.5cm}>{\raggedright\arraybackslash}p{2.9cm}>{\raggedright\arraybackslash}p{4.1cm}>{\raggedright\arraybackslash}p{1.1cm}>{\raggedright\arraybackslash}p{4.4cm}@{}}
    \toprule
    & Invariant & Statement & Holds & What breaks without it \\
    \midrule
    I1 & \textbf{Form invariance}
       & $x \sim_c x'$ implies equivalent decision-relevant outcome
       & No
       & Tests do not generalise past the wording tested \\
    \addlinespace
    I2 & \textbf{Closure}
       & $\Out_S(p)$ and $\Out_S(p \rightarrow q)$ imply $\Out_S(q)$
       & No
       & Conclusions the system's own premises entail are unavailable \\
    \addlinespace
    I3 & \textbf{Decomposition}
       & $\Out_S(p \wedge q)$ implies $\Out_S(p)$
       & No
       & Summaries assert what their parts deny \\
    \addlinespace
    I4 & \textbf{Consistency}
       & Not both $\Out_S(p)$ and $\Out_S(\neg p)$
       & No
       & Two users, or one user twice, get opposite answers \\
    \addlinespace
    I5 & \textbf{Self-report fidelity}
       & Claims about what it knows track what it asserts
       & No
       & Confidence cannot be used to route or gate \\
    \addlinespace
    I6 & \textbf{Correctness}
       & $\Out_S(p)$ implies $p$
       & No
       & Output cannot be relied on unchecked \\
    \addlinespace
    I7 & \textbf{Repeatability}
       & Same input, same session, same assertion
       & No
       & A verified answer is not a verified behaviour \\
    \bottomrule
  \end{tabular}
  \caption{Behavioural invariants. I6 is the one everybody already knows fails, and
    it is the least interesting: nobody deploys these systems believing output is
    true by construction. I1 and I7 are the ones that invalidate validation.}
  \label{tab:invariants-behaviour}
\end{table}

One reading to head off first. These are invariants, so \emph{holds} is a universal
claim and a single counterexample defeats it. ``No'' in the third column therefore
means the property is not guaranteed, not that behaviour is always inconsistent. In
practice these systems satisfy I1--I4 on most inputs most of the time, which is
precisely why the failures are hard to catch in testing and why a specification cannot
rely on them. An invariant that holds usually is not an invariant; it is a tendency,
and nothing downstream can be built on it.

Three further observations, since the pattern matters more than the entries.

\paragraph{I6 failing is priced in; I1 failing is not.} Every deployment I have
seen assumes the system is sometimes wrong and builds review around it. Almost none
assumes that a reviewed, approved, tested behaviour will not reproduce under a
rephrasing. The first is a known unknown that review handles. The second
invalidates the review.

\paragraph{I2--I4 share one cause.} These are not three independent defects. There
is no component whose function is to maintain the set of asserted claims---nothing
that closes it under consequence, decomposes conjunctions, or checks it for
contradiction. By Section~\ref{sec:what-computed} the asserted ``set'' is not a
maintained data structure at all; it is whatever successive forward passes emit.
Absent such a component, I2--I4 have nothing to be true of.

\paragraph{I5 failing has a specific and costly consequence.} A great deal of
proposed architecture routes on model-reported confidence: escalate when the model
says it is unsure, proceed when it says it is confident. I5's failure means the
routing signal is produced by the same process as the content and is not a
measurement of it. Verbalised confidence can be trained to better aggregate
calibration \citep{lin2022}, and models have some capacity to evaluate their own
outputs \citep{kadavath2022}; neither delivers the per-claim reliability a gate
requires. Aggregate calibration and per-claim boundary recognition are different
properties, and only the second supports a decision about \emph{this} answer.

That last point generalises, and it is the shape of the whole argument: a property
that holds on average over a distribution is not a property you can gate on for an
individual case. Accuracy is such a property, as
Proposition~\ref{prop:independence} shows. So is calibration.

\subsection{System invariants: what must be supplied}

The constructive claim. These are conditions on $S$, not on $M$. None holds by
default; each is achievable by architecture; and together they are what must be
true for a response to cross from ordinary generative assistance into consequential
decision support.

\begin{table}[ht]
  \centering
  \small
  \begin{tabular}{@{}>{\raggedright\arraybackslash}p{0.7cm}>{\raggedright\arraybackslash}p{3.3cm}>{\raggedright\arraybackslash}p{4.6cm}>{\raggedright\arraybackslash}p{4.4cm}@{}}
    \toprule
    & Invariant & Requirement & Control that supplies it \\
    \midrule
    I8 & \textbf{Source identity}
       & Every cited source has stable identity and established authenticity
       & Controlled corpus, ingestion authentication \\
    \addlinespace
    I9 & \textbf{Temporal validity}
       & Source version and effective date are recorded and compared against $t$
       & Version metadata, corpus release pinning \\
    \addlinespace
    I10 & \textbf{Context binding}
       & $c$ is explicit before assertion, or explicitly marked unknown
       & Intake resolution, typed request objects \\
    \addlinespace
    I11 & \textbf{Applicability}
       & Evidence is filtered on scope, not similarity alone
       & Metadata-aware retrieval, scope predicates \\
    \addlinespace
    I12 & \textbf{Traceability}
       & Each consequential claim links to the evidence supporting it
       & Claim decomposition, claim-level citation \\
    \addlinespace
    I13 & \textbf{Support verification}
       & The cited evidence is checked to support the claim at the strength asserted
       & Entailment checks, policy logic, human review by tier \\
    \addlinespace
    I14 & \textbf{Reproducibility}
       & The assertion can be re-derived from recorded inputs and versions
       & Snapshot pinning, deterministic replay \\
    \addlinespace
    I15 & \textbf{Authorization}
       & Reliance is permitted at the applicable tier, with duties separated
       & Tiered policy, approval gates, access control \\
    \addlinespace
    I16 & \textbf{Reconstructability}
       & The full decision path is retained for the required period
       & Audit record, retention policy \\
    \addlinespace
    I17 & \textbf{Controlled failure}
       & Unmet conditions produce qualification, refusal, or escalation---not an
         unqualified answer
       & Output gating outside $M$ \\
    \addlinespace
    I18 & \textbf{Sufficiency}
       & The evidence set is checked for coverage of every condition the claim requires
       & Required-field or checklist verification against the rule \\
    \addlinespace
    I19 & \textbf{Defeater search}
       & A scoped search for contrary or undermining authority of equal or higher
         standing is performed and recorded
       & Governed defeater query; institution's precedence order \\
    \bottomrule
  \end{tabular}
  \caption{System invariants. I8--I9 are ordinary data governance. I10--I11 are the
    context work that Section~\ref{sec:cell-v} shows is most often absent.
    I12--I13 are what distinguishes grounding from citation. I14--I16 are what an
    audit actually asks for. I17 is what makes the others enforceable rather than
    advisory. I18--I19 correspond to the coverage and no-defeater conditions of
    Definition~\ref{def:grounding}, and exist so that every condition in the definition
    has a system obligation that establishes it.}
  \label{tab:invariants-system}
\end{table}

Two points about the second table.

\paragraph{I17 is the load-bearing one.} Without a controlled-failure path, every
other system invariant degrades into a preference. A system that detects missing
applicability and answers anyway has not enforced I11; it has measured it. And by
Section~\ref{sec:missing} the gate cannot live inside $M$, because generated
refusal is a continuation rather than a branch. I17 is therefore an architectural
requirement, not a prompt.

\paragraph{Proportionality, not maximalism.} Not every use requires all ten.
Section~\ref{sec:tiers} makes the controls proportional to the consequence of
reliance, and a drafting assistant needs none of I13--I19. What does not follow is
the converse: a system cannot claim auditability because it reports an accuracy
score, emits citations, or carries a generic disclaimer. Those are compatible with
failing every invariant in Table~\ref{tab:invariants-system}.


\section{Grounding Is Auditable; Accuracy Is Not}
\label{sec:auditable}

\subsection{The asymmetry}

Section~\ref{sec:invariants} said what does not hold. This section says which property
can be measured instead, and why that property rather than the obvious one.

Accuracy and grounding are both properties of assertions, and they differ in a way
that decides which one a governance regime can use.

To check whether an assertion is \emph{correct}, you need to know the answer.
That means reference answers, which means the check can only be performed offline, on
questions whose answers you already have---which is to say, on questions nobody
needed to ask. You cannot check correctness on the query a user just sent, because
if you could answer it you would not have sent it to the system.

To check whether an assertion is \emph{grounded}, you do not need to know the
answer. You need to know what evidence the system used and whether that evidence
is applicable at $c$. Both are properties of the system's own execution. Condition
(G1) is a check on the retrieved evidence: is it authentic, current at $t$, of
the right jurisdiction and use? Condition (G2) is a perturbation test: ablate or
contradict the evidence and see whether the assertion survives. Neither consults
a reference answer.

This has a consequence out of proportion to how obvious it sounds once stated.
Grounding can be checked \emph{at runtime, per assertion, on live traffic}.
Accuracy cannot. That is the difference between a metric and a gate. A metric
tells you, after the fact, how a sample went. A gate decides whether this
particular assertion is allowed out. Everything a regulated deployment actually
needs---blocking, escalating, qualifying, logging a reconstructable basis---requires
a gate, and only grounding supplies one.

Grounding also does not inherit the I1 problem. Form-invariance failure means that
whether the system asserts $p$ depends on the wording. It does not follow that
whether an assertion \emph{is grounded} depends on the wording: given an assertion,
the question of whether applicable evidence supports it and whether the mechanism
was sensitive to that evidence is settled by the execution trace, not by the
phrasing that elicited it. Form-invariance failure changes \emph{what} gets asserted,
and the grounding check applies to whatever was in fact asserted. So the pathology
of Section~\ref{sec:canonical-form} does not undermine the measurement---it shows
up, correctly, as a varying assertion set with a measurable grounding rate.

\subsection{Determinism is not the requirement}
\label{sec:determinism}

Before the measurement itself, one reading has to be blocked. The asymmetry above may
suggest that auditable systems must be deterministic, and grounding auditable because
generation is not. That reading is false and it would misdirect the remedy.

Plenty of non-deterministic computation is relied upon in regulated settings and
audited without difficulty. Monte Carlo valuation, randomised algorithms, and
sampling-based statistical procedures all produce different output on different
runs, and all are defensible. What makes them so is not determinism but three other
properties: the run is \emph{reproducible}, because the seed and the state are
recorded and the computation can be replayed; the error is \emph{bounded}, because a
convergence or confidence statement is available; and there is a
\emph{specification} saying what the computation was supposed to approximate.

Generation can satisfy the first of those, but only under conditions worth
spelling out, because they are frequently not the conditions of a commercial
deployment. Greedy decoding removes sampling variance. It does not by itself produce
identical output run to run. Floating-point addition is not associative; the order in
which a reduction is performed on an accelerator depends on which kernel is selected;
and kernel selection depends on batch shape. A request batched with different
neighbours can therefore yield slightly different logits, and where two candidate
tokens are close, a different token. Sparse expert routing under capacity limits
introduces the same batch dependence more directly. Behind a hosted endpoint, the
served model version may change without notice.

Reproducibility is therefore an architectural achievement rather than a property of
the technology. It requires pinned model versions, controlled batching or
single-request inference, deterministic kernel selection, recorded decoding
parameters, and a stateless request. Those conditions are available to an
organisation that self-hosts or contracts for them, and largely unavailable to one
consuming a shared public endpoint. This is why I14 appears in
Table~\ref{tab:invariants-system} as something a system must supply rather than
something it has, and a validation programme that assumes replay works should test
that assumption rather than inherit it.

The point for this section is that non-determinism is not the \emph{irreducible}
gap. It can be engineered away, at a cost that is ordinary architecture.

The gap is that reproducibility over a single input buys very little when the set of
inputs you need to reason about is a semantic equivalence class you cannot
enumerate. Replaying one string tells you what that string does. It says nothing
about the materially equivalent string a different user will send tomorrow, and by
Section~\ref{sec:canonical-form} there is no bound to appeal to in its place. The
contrast with adversarial robustness is instructive: there, perturbations live in a
metric space, so a claim of the form ``behaviour is stable within $\varepsilon$'' is
both meaningful and sometimes certifiable. Paraphrase admits distance measures---
embedding cosine, edit distance, entailment scores---but none with the property the
certification argument needs: that everything within $\varepsilon$ of a request is
materially equivalent to it, and everything materially equivalent is within
$\varepsilon$. Semantically identical requests can be far apart in embedding space
and materially different requests can be adjacent, which is why the corresponding
guarantee cannot be stated even approximately.

So the requirement is reproducibility plus a characterisation of behaviour over the
class of inputs that matter. The first is engineering and this paper treats it as
such. The second is the one with no known general solution, which is why
Section~\ref{sec:auditable} argues for measuring grounding---a property of the
individual assertion and its evidence---rather than attempting to certify behaviour
over a class that cannot be enumerated.

The distinguishing property is worth stating once, because the whole argument turns
on it: the systems in scope produce, as their primary output, statements that people
treat as claims about the world, and they produce them by sampling token sequences.
The first half is why the question of grounding arises at all. The second half,
developed in Section~\ref{sec:what-computed}, is why it cannot be answered by
inspecting the model.

Two boundary cases. Multimodal systems are in scope to the extent that their
text output is produced as in (1)--(4), and out of scope for their image or audio
output. Retrieval-augmented and tool-using systems built on models in scope are
squarely in scope---indeed they are the systems this paper is chiefly about, since
Section~\ref{sec:measurability} argues that only such systems can be measured at
all.

Later developments---longer contexts, sparse routing, tool use, and
reinforcement training on verifiable outcomes---change capability substantially and
change none of the structural properties identified below, a claim I return to in
Section~\ref{sec:better-models}.

\subsection{What independence implies for a validation programme}

Proposition~\ref{prop:independence} established that accuracy and the grounding
rate vary independently over the whole feasible region. Read as a statement about
validation programmes rather than about models, it says something sharper than
``accuracy is incomplete.'' It says accuracy is \emph{uninformative} about the
property that governs reliance: a reported accuracy figure is consistent with every
grounding rate, so no accuracy threshold, however high, constitutes evidence that a
system asserts within its evidence.

The asymmetry above explains why this is not symmetrical bad news. Both properties
are worth knowing, and only one of them can be established for the assertion in
front of you at the moment you need to decide about it.

\subsection{The measurement target}

The quantity to measure is the rate at which assertion outruns applicable
evidence---which, by Definition~\ref{def:hallucination}, is the hallucination rate:
\begin{equation}
  \text{hallucination rate}
  \;=_{\mathrm{def}}\;
  \Pr\bigl(\Hall(p \mid c) \;\big|\; \Out_S(p)\bigr)
  \;=\;
  \Pr\bigl(\neg\,\Grd_S(p \mid c) \;\big|\; \Out_S(p)\bigr).
  \label{eq:rate}
\end{equation}

This is worth pausing on, because it is where the redefinition earns its keep. Under
the standard truth-conditional definition, a hallucination rate can only be
estimated against reference answers, offline, on questions whose answers are already
known. Under Definition~\ref{def:hallucination} it is estimated from the system's own
execution: what evidence was used, and did it apply. Same word, same intuition, and a
quantity that can be computed on live traffic and used to stop an assertion. The
redefinition is what converts a retrospective score into an operable control.
This is \eqref{eq:hall} turned into a statistic. It is estimated by attribution
audit rather than comparison against reference answers, and the framework for attributing
generated statements to identified sources is established
\citep{rashkin2023}: sample assertions, and for each determine whether it is
attributable to identified evidence that is applicable at $c$.

Report it decomposed, not pooled. The three rows of
Table~\ref{tab:remedies} have different owners, so a single figure defeats the
purpose of having measured anything:
\begin{itemize}
  \item ungrounded rate---no evidential connection (cells ii, iv);
  \item misapplied rate---evidence present but inapplicable at $c$ (cell v);
  \item error rate---grounded and false (cell iii), reported against the corpus
        rather than the system.
\end{itemize}

\subsection{Where this can be measured, and where it cannot}
\label{sec:measurability}

One constraint is essential, and it is the point at which this paper's argument
becomes an argument about architecture rather than about metrics.

The evidential base has two parts and they are not equally tractable.

For \textbf{retrieval-provided evidence}, condition (G2) is directly testable.
The evidence is identified, it is present in the context, and you can ablate it,
contradict it, or perturb it and observe whether the assertion survives. The test
is cheap, automatable, and runs on live traffic.

For \textbf{parametric evidence}---whatever the training corpus contributed to the
weights---it is not. ``Had $E$ not supported $p$'' ranges over counterfactual
training corpora, and evaluating it properly means retraining. Training-data
attribution methods approximate the counterfactual, and at frontier scale they
remain a research problem rather than an instrument. The difference between the two
cases is one of kind, not degree.

So \eqref{eq:rate} is well defined for any system and estimable only for systems
that make their evidential base explicit at inference time. An unaugmented model's
ungrounded assertion rate is a perfectly good quantity that nobody can currently
measure.

This is not a limitation of the definition. It is a requirement on the system, and
it is the strongest practical conclusion of this paper. If the rate at which a
system asserts beyond its applicable evidence is the quantity that governs whether
you may rely on it, then the system has to be built so that the quantity is
observable: an identified, versioned, context-resolved evidential base at the
point of assertion, not a training corpus somewhere behind it.
Section~\ref{sec:architecture} sets out what that requires.


\section{Governed Retrieval: A Reference Architecture}
\label{sec:governed-rag}

Section~\ref{sec:measurability} established a requirement rather than a preference:
the grounding rate is measurable only where the evidential base is explicit at
inference time. Section~\ref{sec:retrieval-insufficient} established that standard
retrieval supplies one of seven grounding dimensions reliably and none of the three
that are relations to $c$. This section sets out what closes the gap.

The distinction I will use throughout:

\begin{quote}
Standard retrieval retrieves evidence. \emph{Governed} retrieval establishes the
controlled conditions under which retrieved evidence may support an institutional
assertion.
\end{quote}

This is architecture, not prompt engineering. None of the eight components below can
be implemented as an instruction to the model, and Section~\ref{sec:missing} says
why: instructions enter as tokens and are processed by the same mechanism as the
content, so an instruction to abstain is a probabilistic influence on a
continuation rather than a precondition on output.

\subsection{The Holodeck principle}
\label{sec:holodeck}

Before the components, the requirement that motivates the first two of them.

\begin{quote}
Natural language is an incomplete specification for consequential action.
\end{quote}

A fluent request routinely leaves unresolved: the referents; the jurisdiction; the
effective time; the applicable policy; the objective; the requester's authority; the
constraints; the side effects; the reversibility; the threshold at which action
should be taken; and whether the requester wants analysis, a recommendation,
execution, or a binding determination.

The point is not that natural language is a poor interface. It is an excellent
interface, which is why these systems are worth building. The point is that a
fluent request cannot be treated as a \emph{complete operational specification} in a
regulated decision or execution workflow, and that the gap between the two is
invisible in the request itself. A well-formed sentence does not announce which of
its parameters were never bound.

Two consequences. The system must force or validate material bindings before
crossing into consequential action---which is I10. And the boundary at which it
crosses has to be an explicit architectural decision, because by
Section~\ref{sec:rlhf} the training regime supplies pressure to answer rather than
to ask.

\subsection{Eight components}

\paragraph{1. Authoritative corpus control.} Sources are approved, authenticated,
and governed. Each carries stable identity, version, effective date, jurisdiction,
authority level, and applicability metadata, and sources can be pinned to a corpus
release or a decision-time snapshot. This supplies I8, I9, and the substrate for
I11 and I14. It is the component most familiar from ordinary data governance and,
in my experience, the one most often assumed to be sufficient on its own.

\paragraph{2. Context resolution before assertion.} Material variables---$j$, $t$,
$f$, $u$---are resolved, collected, or explicitly marked unknown before an
assertion is produced. Where material context is missing the system requests it,
qualifies the output, declines to produce a decision-ready assertion, or escalates.
This supplies I10 and is the only defence against illicit universalization
(Section~\ref{sec:illicit}). It is also the component with the highest product cost,
because it makes the interaction less immediately satisfying, and it is therefore
the first to be cut.

\paragraph{3. Metadata-aware retrieval.} Retrieval filters or ranks on scope
constraints, not on semantic similarity alone: jurisdiction predicates, effective
date windows, population or product class, decision type. The retrieval event and
the resulting evidence set are retained. This supplies I11, and it is the direct
answer to the observation of Section~\ref{sec:retrieval-insufficient} that a
paediatric and an adult dosing passage are near-neighbours in embedding space.

\paragraph{4. Claim decomposition.} Complex outputs are decomposed, where feasible,
into atomic claims, each with linked evidence and an explicit support relation. This
supplies I12. A response-level citation list does not permit checking the third
sentence, and the third sentence is frequently the one relied upon.

\paragraph{5. Evidence-to-assertion verification.} Citation presence is
insufficient. The system checks whether the cited evidence supports the claim, and
whether it does so under $c$ and at the strength asserted
(Section~\ref{sec:entails-supports}). Controls may combine deterministic rules,
structured policy logic, source checks, independent model review, and human
validation, proportioned to risk. This supplies I13 and is where the
$\Sensitive$ condition of Definition~\ref{def:grounding} is actually enforced
rather than assumed.

\paragraph{6. Output gating.} If evidence is missing, stale, conflicting,
inapplicable, or insufficient for the strength of claim requested, the system does
not present an unqualified assertion as decision-ready. It returns bounded
information, an explicit uncertainty, a request for missing context, or an
escalation path. This supplies I17, and it must sit outside $M$.

\paragraph{7. Audit record and replay.} Retain the prompt or structured input; the
resolved context; source identifiers and versions; retrieved passages; model
version; prompt or template version; system configuration; verification results;
human interventions and approvals; the final response; and timestamps with
retention identifiers. This supplies I14 and I16, and it is what answers the
question posed at the start of Section~\ref{sec:intro}.

\paragraph{8. Authorization and tiered reliance.} Permitted use depends on the
consequence of reliance, not only on model quality. Higher-consequence uses require
stronger constraints, independent validation, and where appropriate a qualified
human decision-maker with the standing to be accountable. This supplies I15 and is
developed in Section~\ref{sec:tiers}.

\subsection{The flow}

\begin{center}
raw generative output
$\rightarrow$ context resolution
$\rightarrow$ controlled retrieval
$\rightarrow$ applicability filtering \\
$\rightarrow$ evidence-to-claim verification
$\rightarrow$ authorization or escalation
$\rightarrow$ auditable decision support
\end{center}

Two things about this sequence carry most of its value, and neither is the ordering.

Each arrow is a \emph{stopping point}. Context resolution can return a request for
missing bindings; applicability filtering can exclude every candidate source; verification
can find that the evidence does not reach the claim; the defeater search can find contrary
authority; and authorization can decline. A pipeline in which the only outcome is an
answer has implemented the boxes without the arrows, and it is the arrows that make the
conditions enforceable rather than advisory.

Each arrow is also a \emph{record}. What the audit trail contains is not a summary of the
final answer but the sequence of determinations that produced it: which context was
resolved and how, which sources survived filtering and which were excluded and why, what
the verification returned, what the defeater search returned, and who authorised the
result. That is what makes the flow reconstructable months later by someone who was not
present, which is the requirement Section~\ref{sec:intro} opened with.

The sequence is not a confidence gradient. It is a sequence of distinct conditions, and an
assertion that satisfies six of them is not ninety per cent warranted; it is unwarranted
for a reason that the record should name.

\subsection{What each intervention buys}

Common interventions address different invariants, and conflating them is how programmes
spend a great deal and move one number. The table below is best read as a purchasing
guide: the left column is what a programme funds, the middle column is what that
purchase actually delivers against Table~\ref{tab:invariants-system}, and the right
column is what it is frequently assumed to deliver and does not.

Two features of it are worth pointing out before reading it. Retrieval appears with two
distinct entries in the middle column because it does two unrelated things: it reduces
fabrication, which is why it is usually funded, and it supplies an identified evidential
base, which is what makes \eqref{eq:rate} measurable at all and is the more durable of
the two. And no single row delivers more than a handful of invariants, which is the
practical form of the argument in Section~\ref{sec:three-teams}: there is no purchase
that makes a system auditable, only a set of purchases that between them cover the
conditions.

\begin{table}[ht]
  \centering
  \small
  \begin{tabular}{@{}>{\raggedright\arraybackslash}p{3.3cm}>{\raggedright\arraybackslash}p{4.6cm}>{\raggedright\arraybackslash}p{5.4cm}@{}}
    \toprule
    Intervention & Buys & Does not buy \\
    \midrule
    Standard retrieval
      & I8; makes \eqref{eq:rate} measurable; reduces \textsc{Fabrication}
      & I10--I11, I13; form invariance; applicability \\
    \addlinespace
    Better base model
      & fewer fabrications on covered topics
      & I8--I19; may worsen \textsc{Lucky} by guessing better \\
    \addlinespace
    Verifier loops
      & I13 on checkable claims; some of I2--I4
      & grounding, unless the verifier checks evidence rather than form \\
    \addlinespace
    Persistent memory
      & a precondition for I7 and I14
      & either of them---storage is not retrieval discipline \\
    \addlinespace
    Context resolution
      & I10, and the substrate for I11
      & anything on the generation path \\
    \addlinespace
    Output gating
      & I17, and thereby enforceability of the rest
      & correctness of the evidence itself \\
    \bottomrule
  \end{tabular}
  \caption{Retrieval does double duty in most deployments: it is funded to improve
    accuracy, and its more durable contribution is making the grounding rate
    observable at all.}
  \label{tab:interventions}
\end{table}

The row worth dwelling on is the second. A better base model improves the numbers
most programmes report while leaving every system invariant untouched, and its
effect on the lucky cell runs the wrong way: a model that guesses better produces
more cell-(ii) assertions and scores higher for doing so.

\subsection{Structured evidence is the strongest case}
\label{sec:structured}

The components above are described in terms of retrieved passages, which
understates what is available. The highest-assurance form of $E$ in an enterprise is
not a document at all. It is a \emph{query against a system of record}, with the query
text, the parameters, the result set, and the as-of state all recorded.

Structured evidence satisfies the grounding conditions more cleanly than text
retrieval does, because the evidence has a schema. Authenticity follows from the source
system's own controls. Versioning is the table's temporal model rather than a document
revision. Applicability becomes a set of predicates over columns---jurisdiction,
effective date range, product class, population---evaluated rather than judged.
Sufficiency becomes a completeness check over required fields. Traceability is the
recorded query. And the defeater check of \eqref{eq:nodefeater} becomes a query for
contradicting rows rather than a search for contradicting prose.

The practical consequence is a preference order that most retrieval-first designs
invert: where a claim can be answered from a governed table, it should be, and the
model's role is to render the result rather than to find it. Text retrieval is the
fallback for claims whose evidence exists only as prose, and it inherits every
judgement that the schema would otherwise have decided.

\subsection{What this costs}
\label{sec:cost}

The controls are not free and a proposal that does not say so is not usable. The
multipliers can be given structurally, without benchmarks, because they follow from how
many model calls and retrievals each control requires.

Let a response contain $m$ consequential claims, each with up to $k$ candidate
supporting sources. Generation is one call. Claim decomposition is one call plus $O(m)$.
Support verification is $O(mk)$ calls where it is model-based, and $O(1)$ per claim where
the check is deterministic. The defeater search is one additional retrieval per claim,
$O(m)$. Coverage checking is $O(m)$ against an enumerated rule and unbounded against a
rule that must be read out of prose. So the dominant term is $mk$, and the design levers
follow directly: reduce claims per response, reduce candidate sources per claim, or
replace model verification with deterministic checks. A ten-claim response over three
candidate sources each is on the order of thirty verification calls against one
generation call---which is the difference between a system that costs what a chat costs
and one that does not.

Storage grows with retained evidence rather than with traffic, and
Section~\ref{sec:retention} constrains what may lawfully be kept.

\begin{table}[ht]
  \centering
  \small
  \begin{tabular}{@{}>{\raggedright\arraybackslash}p{3.9cm}>{\raggedright\arraybackslash}p{3.1cm}>{\raggedright\arraybackslash}p{6.6cm}@{}}
    \toprule
    Component & Cost profile & Notes \\
    \midrule
    Corpus control, metadata
      & one-off plus upkeep
      & Ontology and curation effort; no per-request cost \\
    \addlinespace
    Metadata-aware retrieval
      & negligible per request
      & Predicate filtering is cheaper than the embedding search it constrains \\
    \addlinespace
    Context resolution
      & one round trip, or a form
      & Latency and interaction cost, not compute \\
    \addlinespace
    Claim decomposition
      & multiplies calls
      & Linear in claims per response \\
    \addlinespace
    Support verification
      & multiplies again
      & The dominant cost. Deterministic checks are cheap; model-based
        entailment is per claim per source; human adjudication is the expensive tail \\
    \addlinespace
    Defeater search (I19)
      & one extra retrieval
      & The governed contrary query, run under the same applicability predicates.
        Cheap, and it needs the precedence order to exist \\
    \addlinespace
    Coverage checking (I18)
      & per claim
      & Cheap where the rule's conditions are enumerated; expensive where they must be
        inferred from prose \\
    \addlinespace
    Audit record
      & storage, growing
      & See Section~\ref{sec:retention} on what may lawfully be kept \\
    \bottomrule
  \end{tabular}
  \caption{Cost profile of the controls. The expensive items are per-claim
    verification and human adjudication, which is precisely why
    Section~\ref{sec:tiers} makes them proportional to consequence.}
  \label{tab:cost}
\end{table}

The tiering of Section~\ref{sec:tiers} exists for this reason and not only for
governance tidiness. If every request bore the full control set the economics would
not work, and a design that applies Tier~3 controls to Tier~1 traffic will be
abandoned rather than descoped.

\subsection{Retention against minimisation}
\label{sec:retention}

The audit record of component~7 conflicts with another control regime, and the
conflict should be stated rather than discovered in review. Prompts, resolved contexts,
and retrieved passages routinely contain personal or special-category data. Data
minimisation and erasure obligations cut directly against the retention that
reconstructability requires.

Four resolutions, none complete on its own. Retain references, identifiers, and hashes
in place of content wherever the content is recoverable from a governed source, so the
audit record points at evidence rather than copying it. Segregate the audit store and
control access to it separately from the application. Align retention to the retention
schedule of the \emph{decision} the assertion supported, rather than choosing a new
period for the AI system. And where erasure is exercised, preserve the structure of the
record---that a claim was made, on what evidence class, under what authority---while
removing the identifying content.

The point for this paper is narrower than the resolutions. Reconstructability is a
requirement in tension with minimisation, the tension is a design decision rather than
an implementation detail, and it requires a fifth function: privacy and records
management, which Section~\ref{sec:three-teams} adds to the three that own the rest.

\subsection{What this does not claim}

Governed retrieval does not make these systems infallible, and nothing above should
be read as proposing that it does. Cell (iii) survives it entirely: a controlled,
versioned, applicable, authoritative corpus can still be wrong, and no amount of
architecture converts a bad source into a good one. Verification is relative to the
checks implemented and the assumptions they encode. Human review introduces its own
failure modes and its own capacity limits. And the floor of
Section~\ref{sec:floor} does not move.

The function of the architecture is narrower and worth stating precisely. It
transforms otherwise unlicensed generated text into a controlled, inspectable input
to an accountable institutional process---one whose failures are attributable,
whose evidence is identifiable, and whose reasoning can be reconstructed after the
fact by someone who was not present. That is a lower bar than reliability and a
much higher bar than accuracy.


\section{Bounding the Request Space}
\label{sec:bounded}

Section~\ref{sec:canonical-form} argued that the absence of a canonical form is what
makes accuracy uncertifiable, and Section~\ref{sec:normaliser} argued that training a
semantic normaliser is not the remedy. There is a third option, and it is the one an
architect reaches for first: if the model will not canonicalise its input and cannot
be taught to, \emph{do not give it free-form input at all}.

Restricting the interface is not a workaround. It is the direct application of the
diagnosis: canonicalisation moves upstream of the model, into a schema, where it is
inspectable and versioned. This section works through what that buys, what it does
not, and where the residual risk goes---because it does not disappear.

\subsection{The construction}

Replace free text with a typed request. Instead of ``is this approved for children?''
the interface accepts an object:

\begin{center}
\small
\begin{tabular}{@{}l>{\raggedright\arraybackslash}p{7.2cm}@{}}
\toprule
\texttt{intent} & \texttt{INDICATION\_CHECK} (from a closed set) \\
\texttt{subject\_id} & a key into a system of record \\
\texttt{indication} & a coded value \\
\texttt{jurisdiction} & required \\
\texttt{as\_of} & required \\
\texttt{population} & required, from a controlled vocabulary \\
\bottomrule
\end{tabular}
\end{center}

The object is rendered into a fixed template by code the organisation controls and
versions. Requests are stateless: one request, one fresh context.

The fields are not arbitrary, and the discipline that selects them is the point rather
than the example. Every component of $c$ appears, because
Section~\ref{sec:resolved-context} established that a claim whose context is unresolved is
not yet evaluable---so jurisdiction, valid time, and intended use are required fields, and
a request missing any of them does not validate. The material facts $f$ enter as keys into
systems of record rather than as free text, so that the entity under discussion is
identified rather than described. The intent is drawn from a closed set, because the
decision class determines which sources have standing under $\standing$ and therefore what
counts as authoritative. And coded values replace prose wherever a controlled vocabulary
exists, since that is what allows applicability to be evaluated as a predicate in
Section~\ref{sec:governed-rag}'s third component rather than judged.

What the schema encodes, in other words, is the set of bindings that would otherwise have
to be guessed. Anything the type system requires is a thing the system can no longer
silently assume.

\subsection{What this fixes}

Against Table~\ref{tab:invariants-system}, the gains are specific and substantial.

\paragraph{I1 becomes satisfied by construction.} Not because the model acquired form
invariance, but because the opportunity to violate it has been removed. Two
materially equivalent requests are now \emph{the same request}: the equivalence class
of Section~\ref{sec:canonical-form} has exactly one member, and the canonical form is
supplied by the schema.

\paragraph{The request space becomes enumerable.} It is a product of typed fields
rather than $\mathcal{V}^{*}$. It may be combinatorially large, but it is
\emph{stratifiable}: coverage can be designed by jurisdiction, date band, population,
and intent, with known sampling properties. That is ordinary test design. The
validation problem becomes well posed in the sense of
Section~\ref{sec:not-easier}---the same sense in which it is well posed for a
classifier.

\paragraph{I10 is enforced by the type system rather than detected.} A request without
a jurisdiction does not validate. The misapplied category of
Section~\ref{sec:cell-v} does not need to be caught downstream, because the binding
it depends on is a precondition of submission. This is the single largest gain, and
it is available with no machine learning at all.

\paragraph{I7 becomes achievable.} Statelessness removes context accumulation, which
is the dominant source of within-session drift. Combined with the reproducibility
conditions of Section~\ref{sec:determinism}, replay becomes a realistic control
rather than an aspiration.

\subsection{What this does not fix}

\paragraph{The output space is still $\mathcal{V}^{*}$.} Bounding the input bounds the
input. The model can still emit any string, so ungrounded assertion survives
untouched: a confident figure resting on nothing is exactly as available as before.
Grounding is I13 and no amount of input discipline supplies it.

\paragraph{Errors are still not label mismatches.} If the response is prose, it can be
partly right, right for the wrong reason, or adequate for a weaker claim than the one
made. The atomic countable error of the classifier case does not follow from a typed
input.

\subsection{Where the problem goes}
\label{sec:relocation}

This is the part that decides whether the design is an improvement or a disguise.
\textbf{Who fills in the request?}

If a person translates their need into the typed object, the ambiguity has not been
eliminated; it has been moved into that translation, and it is unrecorded unless the
system captures it.

If a language model performs the translation---which is what is usually built,
because it preserves the interface that made the technology attractive---then the
unbounded input space has been reintroduced at the parsing stage. Form-invariance
failure now determines \emph{which request object is constructed}. A paraphrase can
produce a different jurisdiction, a different population, or a different intent, and
everything downstream will process that object correctly and produce a well-grounded
answer to the wrong question.

This is worse than the original failure in one specific respect: it is less visible.
An ungrounded assertion at least looks like an assertion, and the audit record shows
what evidence was and was not used. A misparsed request produces a record that is
internally consistent and complete. Every control passes. The evidence applies to the
context in the record, and the context in the record is not the user's context.

The mitigation is not technical. It is to \textbf{render the resolved request back and
require confirmation before proceeding}, which converts a silent failure into a
visible one and puts the binding decision where accountability already sits. This is
the pattern enterprise systems have always used before consequential execution, and
the argument of this paper is that a natural-language front end does not earn an
exemption from it. Where confirmation is impractical, the parse should be treated as
part of the assertion for audit purposes: recorded, versioned, and measured for
accuracy against reviewed samples like any other component.

\subsection{Constraining the output, and what that does and does not buy}
\label{sec:constrained-output}

The natural complement is to constrain the output as well: grammar-constrained
decoding into a schema, with an enumerated verdict field, required evidence
references, and an explicit \texttt{INSUFFICIENT\_EVIDENCE} variant.

This is worth doing and its benefits should be stated exactly, because it is easy to
credit it with more than it delivers.

It \emph{does} guarantee that the output parses, that the output space is enumerable,
and that errors become atomic and countable again---restoring on the output side the
well-posedness that typed requests restore on the input side. It \emph{does} make
abstention a first-class value rather than a hoped-for phrasing, and it makes
per-field evidence linkage (I12) structurally enforceable rather than conventional.

It does \emph{not} satisfy I17. A grammar constrains which strings can be emitted; it
does not determine which the model selects. The choice between a substantive verdict
and the abstention variant is made by the same next-token computation as everything
else, subject to the same pressures identified in Section~\ref{sec:rlhf}. Constrained decoding makes abstention \emph{expressible}, not
\emph{correctly triggered}.

I17 is therefore satisfied by the check, not by the schema. The evidence conditions of
Definition~\ref{def:grounded-evidence} are evaluated outside the model, and the gate
must be able to \emph{override} a substantive verdict the model produced---replacing it
with a qualification, a refusal, or an escalation---when those conditions are not met.
A system that trusts the model's own choice of variant has implemented a vocabulary
for abstention and no gate.

\subsection{Where this is practical}

The design suits high-volume, repetitive, well-typed decisions: eligibility
determinations, coding against a rulebook, standard disclosures, reconciliations,
completeness checks. The user is an operator working a queue against a known decision
class, and the intent set is small and stable.

It suits exploratory work badly, and that matters, because exploration is a large part
of why these systems were adopted. An analyst pursuing an unanticipated question has
no intent to select.

Two costs deserve naming.

\paragraph{The intent taxonomy is an ontology project.} Enumerating decision classes,
coding their fields, and maintaining controlled vocabularies is the same work as
building a conformed dimensional model, with the same effort profile and the same
tendency to be underestimated.

\paragraph{An incomplete taxonomy fails quietly.} A question outside the intent set
does not raise an error. It is routed to the nearest available intent, which is a
misparse of the kind described in Section~\ref{sec:relocation}. Twenty intents may
cover most of the volume, and the residual tail is disproportionately where the
unusual, high-consequence cases live. A closed intent set therefore needs an explicit
\emph{no matching intent} outcome and a route for it, or the constraint will silently
manufacture the failure it was introduced to prevent.

\subsection{The effect on what the system is}

One consequence should be faced directly, because it changes how the component should
be described and governed.

If the input is typed, the output is schema-constrained, and the evidence conditions
are checked outside the model, then it is fair to ask what the model is contributing.
Frequently the honest answer is: retrieval, summarisation, and the rendering of a
determination into readable language---while the determination itself is made by
deterministic rules over structured evidence.

That is a legitimate and valuable architecture, and it is one enterprises already know
how to validate. But it \emph{reclassifies} the system. The model is a retrieval and
presentation layer over a decision service, not the decision-maker, and the assurance
argument should say so. A programme that describes such a system as ``AI-driven
decisioning'' is overstating the model's role and, more importantly, misdirecting its
own validation effort toward the component that is doing the least consequential work.

There is a further consequence worth facing, because it follows from combining this
section with Section~\ref{sec:structured} and a reader will combine them. If structured
evidence is the strongest available form of $E$, and typed requests plus constrained
output are what reach the higher reliance tiers, then the architecture that best
satisfies this paper's conditions is one in which the model does comparatively little of
the deciding. Taken to its conclusion, that is an argument for routing decidable claims
\emph{away} from the model altogether.

I take that to be a feature rather than an embarrassment. A framework that identifies
which decisions do not need a language model is more useful than one that assumes every
decision does. What remains for the model, once the decidable residue is routed away, is
the genuinely open-textured part: synthesis across heterogeneous sources, questions whose
evidence exists only as prose, interpretation where the governing text is ambiguous, and
interaction with people who cannot state their need in a schema. That is a smaller claim
about the technology than the field usually makes, and a more defensible one---and it is
the part where the measurement problem of Section~\ref{sec:auditable} genuinely bites,
because it is the part no schema can decide.

The tiering of Section~\ref{sec:tiers} is where this resolves. Bounding the request
space is not a general answer to the problem of ungrounded assertion; it is the
mechanism by which a narrow, well-understood decision class can be moved to a higher
reliance tier than free-form interaction can support. Free-form interfaces are
appropriate at Tiers~0 to~2. Reaching Tier~3 and above is, in practice, an exercise in
bounding what may be asked and what may be answered.


\section{Scale, Scope, and What Is Actually Tunable}
\label{sec:scale}

The architectural choices of Sections~\ref{sec:governed-rag} and~\ref{sec:bounded}
were presented as though the model were fixed. It usually is not. The question an
architect faces early is whether to consume a large general-purpose model as a
service or to operate a smaller model over a defined domain, and the two options are
routinely compared on accuracy. This section argues that accuracy is close to
irrelevant to the comparison, identifies the two variables that do govern it, and
declines a tempting conclusion that does not follow.

\subsection{What scale does, and what it cannot do}

Scale improves accuracy. That is not in dispute and nothing here argues otherwise.
What matters is what an accuracy improvement licenses, and by
Proposition~\ref{prop:independence} it licenses nothing about the grounding rate. The
two quantities are independent across the whole feasible region, so a model that
scores better tells you it scores better.

There is a stronger point available, and it follows from
Section~\ref{sec:pretraining} rather than from any claim about how well these systems
work. The objective \eqref{eq:objective} contains no term that is a function of
truth. Accuracy gains from scale are therefore gains in modelling the training
distribution, which correlates with truth to the degree the corpus is truthful. They
are not convergence toward truth as an optimisation target, because no term in the
loss distinguishes a true continuation from a plausible false one where the two
diverge. Scale sharpens the estimate of the distribution; it does not introduce a
distinction the objective never drew.

This is worth stating carefully, because it is easy to overstate in either
direction. The claim is not that a larger model is not more often right---it is. The
claim is that being more often right on a distribution is a different property from
standing in an auditable relation to evidence, and that no amount of the first
produces the second. The mechanism that would produce grounding is absent from the
architecture, per Section~\ref{sec:missing}, and adding parameters does not add a
mechanism.

\subsection{Reproducibility is a property of topology, not of size}

Invariant I14 is where the comparison is usually miscast. Reproducibility does not
scale, in either direction, because it is not a function of the model at all. It is a
function of who operates the inference.

Section~\ref{sec:determinism} listed the conditions: pinned model versions,
controlled batching or single-request inference, deterministic kernel selection,
recorded decoding parameters, statelessness. Every one of those is a property of the
serving stack. A parameter count appears in none of them.

What is true, and is the defensible core of the common intuition, is that the
\emph{typical consumption pattern} for large models makes those conditions harder to
obtain. A shared multi-tenant endpoint batches your request with other traffic, so
batch composition is outside your control; sparse expert routing under capacity
limits makes that batch composition affect your output; and the served version may
change without notice. None of that follows from the model being large. It follows
from the model being operated by someone else.

The distinction is not pedantic, because it changes what an architect should do about
it. ``Prefer a small model'' is not actionable advice and is sometimes wrong.
``Obtain version pinning, deterministic serving, and stateless requests---by
self-hosting, by contract, or by choosing a deployment that offers them'' is
actionable, and it is available for open-weight models of any size.

\subsection{Corpus scope changes what can be measured}
\label{sec:corpus-scope}

Here the comparison is genuinely asymmetric, and in a way that has nothing to do with
parameter count and everything to do with what the evidential base contains.

Two of this paper's measurements require establishing a \emph{negative}: that the
evidence does not support a claim. The abstention items of
Section~\ref{sec:abstention} need it by construction---an item whose correct
behaviour is silence is only well formed if you can establish that the answer is
absent from what the system can draw on. And the ungrounded assertion rate
\eqref{eq:rate} needs it for the parametric half of $E$, which is why
Section~\ref{sec:measurability} concluded that half is largely unmeasurable.

Negative claims about a bounded, enumerated, versioned corpus are tractable. You can
determine that a controlled document set does not address a question, by search over
an enumerated and versioned index rather than by inference about what a model
absorbed. For a corpus of any size that is work rather than inspection, and it is
tractable in the way any completeness check over a known population is tractable. Negative claims about a web-scale training corpus are not
tractable with current attribution methods, and the gap is one of kind rather than
maturity: the question is what a model did \emph{not} absorb, and attribution
techniques estimate influence of what it did.

A related asymmetry follows from Section~\ref{sec:floor}. The statistical floor on
ungrounded assertion is highest for facts that are rare in the training
distribution---and \emph{rare} is only definable relative to a known corpus. For a
domain corpus you can characterise coverage and therefore estimate where the floor
sits. For a frontier training set you cannot, so the floor is present, consequential,
and unquantified.

So bounding the corpus does not merely improve accuracy on the domain. It changes
which properties of the system are \emph{measurable at all}, which is a different and
more important kind of gain. It is the same move as
Section~\ref{sec:bounded}, applied to the evidence rather than the request:
enumerability is what makes negative statements available, and negative statements
are what auditing requires.

\subsection{What the dials actually buy}

Collecting the choices available to an architect against the invariants they affect:

\begin{table}[ht]
  \centering
  \small
  \begin{tabular}{@{}>{\raggedright\arraybackslash}p{4.0cm}>{\raggedright\arraybackslash}p{2.4cm}>{\raggedright\arraybackslash}p{6.4cm}@{}}
    \toprule
    Choice & Invariants & Effect on auditability \\
    \midrule
    Typed request interface
      & I1, I7, I10
      & Large. Supplies the canonical form and enforces context binding \\
    \addlinespace
    Controlled, versioned corpus
      & I8, I9, I11
      & Large. Precondition for every evidence condition \\
    \addlinespace
    Bounded, enumerable corpus
      & enables I13, I17 testing
      & Large. Makes negative evidence claims and abstention items possible \\
    \addlinespace
    Schema-constrained output plus external gate
      & I12, I13, I17
      & Large. Atomic errors; abstention enforced rather than expressed \\
    \addlinespace
    Operational control of inference
      & I14, I7
      & Moderate to large. Version pinning, batching, statelessness \\
    \addlinespace
    Claim decomposition
      & I12
      & Moderate. Claim-level rather than response-level linkage \\
    \addlinespace
    \textbf{Model parameter count}
      & \textbf{none directly}
      & \textbf{No direct effect.} Affects accuracy, which by
        Prop.~\ref{prop:independence} bounds nothing. One indirect path matters---see
        below \\
    \bottomrule
  \end{tabular}
  \caption{Architectural choices ranked by effect on auditability. Model size is the
    dial most often debated and the one with the least direct effect.}
  \label{tab:dials}
\end{table}

The ordering is the practical conclusion of this section, and it is somewhat
counterintuitive: the variable that dominates procurement discussion has the smallest
\emph{direct} effect on whether the resulting system can be audited.

One indirect path is real and should not be buried in a table. Where an in-scope model
performs the context resolution or request construction of
Section~\ref{sec:relocation}---translating free text into a typed request---its
capability bears directly on I10, and a misparse there produces the most dangerous
failure in the architecture: a complete, internally consistent record answering the
wrong question. So model capability does matter, at exactly one point, and the
conclusion to draw is not that a larger model fixes it. It is that the parse must be
made \emph{reviewable}, per Section~\ref{sec:relocation}, because a control that depends
on model capability is a control whose failure rate is unknown. Better capability
narrows that failure rate without bounding it; confirmation bounds it.

\subsection{The conclusion that does not follow}
\label{sec:not-follow}

It is tempting to conclude that small domain models can be governed and large
general-purpose models cannot. That does not follow, and it is worth saying why,
because the error would misdirect real spending.

Every condition in Table~\ref{tab:invariants-system} is a property of the system
surrounding the model. None mentions parameter count, training corpus size, or
provenance of the weights. A large model placed inside a governed architecture---typed
requests, controlled corpus, external evidence gate, defeater search, audit record,
pinned version---satisfies those conditions exactly as a small model would, because the
conditions are about the architecture. Conversely, a small model consumed through a
shared endpoint with no version pinning and free-form input satisfies almost none of
them.

The two genuine axes identified above are \emph{operational control of the inference
stack} and \emph{enumerability of the evidential base}. Both correlate with model
size in current practice, and neither is caused by it. Self-hosted open weights give
operational control at any scale. A retrieval corpus can be bounded and versioned
whatever model reads it.

The claim I would defend is therefore narrower than the intuition and more useful
than it: \textbf{you cannot govern inference you neither operate nor contract for, and
you cannot make negative evidence claims about a corpus you cannot enumerate.}

The contractual half deserves emphasis, because enterprises govern third-party
processors routinely and the instruments are familiar. What would otherwise be
configuration becomes a term, and the terms are nameable: version pinning with a
defined support window; advance notice of model and serving changes; a deterministic or
single-tenant serving option; recorded decoding parameters; retention of request and
response for a stated period; audit rights over the serving configuration; and
restrictions on sub-processing. An organisation that cannot obtain those has not
discovered a property of large models---it has discovered the terms it accepted. The
practical conclusion is a procurement checklist rather than a prohibition, and it is
the one place in this paper where the binding control is legal rather than technical.

One asymmetry does survive, and it is worth stating because it is where the intuition
is right. Scale is offered as a solution to accuracy, and accuracy is not the
constraint. So an organisation that responds to a grounding problem by buying a larger
model has bought an improvement it cannot use for the purpose it needs, and has
usually given up operational control to obtain it. That is not an argument against
large models. It is an argument against treating model capability as the lever when
the binding constraint is architectural.


\section{Institutional Warrant and Tiered Reliance}
\label{sec:tiers}

The architecture of Section~\ref{sec:governed-rag} serves a condition I have so far
used without defining. This section defines it and makes the controls proportional
to consequence.

\subsection{Warrant}

Grounding is necessary for reliance and not sufficient. An assertion can rest on
authentic, applicable, authoritative evidence, be traceable to it, and still be
something the institution is not entitled to act on---because nobody checked it, or
because nobody with the standing to authorise reliance did so.

\begin{definition}[Institutional warrant]
\label{def:warrant}
\begin{equation}
  \begin{aligned}
  \Warranted_I(p \mid c) \;\leftrightarrow\;
    &\Resolved(c) \wedge \Grd_S(p \mid c)\\
    &\wedge\ \Verified(V, p, c) \wedge \Authorized(I, p, c).
  \end{aligned}
  \label{eq:warrant}
\end{equation}
\end{definition}

where

\begin{itemize}
  \item $\Resolved(c)$: the material context variables of \eqref{eq:context} are
        explicit rather than assumed or inferred without validation;
  \item $\Grd_S(p \mid c)$: Definition~\ref{def:grounding}---which unpacks to the
        six evidence conditions of \eqref{eq:grounded} plus $\Sensitive$, so
        authenticity, versioning, authority, and applicability are preconditions of
        valid support rather than separate requirements;
  \item $\Verified(V, p, c)$: an appropriate procedure has checked the assertion,
        its evidence relation, and the controls the decision class requires;
  \item $\Authorized(I, p, c)$: the institution permits reliance at the relevant
        risk tier, possibly subject to human review or approval.
\end{itemize}

Three properties of \eqref{eq:warrant} matter for how it is used.

It is a conjunction, not a score. There is no aggregation of the four conjuncts into
a confidence figure, and an assertion satisfying three of them is not
three-quarters warranted. This is the same point as the flow in
Section~\ref{sec:governed-rag}: these are distinct conditions with distinct owners
and distinct failure records.

It is a property of a process, not of a model. Nothing in \eqref{eq:warrant} is a
property of model weights, a benchmark score, or a generated response.
$\Authorized$ in particular is a fact about an institution's governing rules and
delegations, and no improvement to $M$ can supply it. This is why warrant cannot be
procured: a vendor can sell you a component that makes grounding measurable, and
cannot sell you an authorization.

It is revocable and context-relative. \eqref{eq:warrant} is not a claim that $p$ is
certain, or even that $p$ is true. It is the condition under which an institution
may responsibly act on $p$ for a defined purpose, knowing that it may be wrong and
that the record will show what was relied upon and why. That is the standard
ordinary institutional decision-making already meets, and the standard machine
assertion currently does not.

\subsection{Reliance tiers}

Controls should be proportional to the consequence of reliance. The table below is
deliberately compact: its purpose is to show that permissible reliance is set by
consequence rather than by model quality, and that argument does not require a
fully specified control catalogue. Any real programme will need its own tiering,
calibrated to its regulator and its risk appetite.

\begin{table}[ht]
  \centering
  \small
  \begin{tabular}{@{}>{\raggedright\arraybackslash}p{1.1cm}>{\raggedright\arraybackslash}p{3.6cm}>{\raggedright\arraybackslash}p{4.4cm}>{\raggedright\arraybackslash}p{4.4cm}@{}}
    \toprule
    Tier & Use & Required & Invariants engaged \\
    \midrule
    0 & Unwarranted generative assistance: brainstorming, drafting
      & Clear framing as unverified; no reliance claim
      & none \\
    \addlinespace
    1 & Informational assistance
      & Visible uncertainty; no reliance claim
      & I17 \\
    \addlinespace
    2 & Evidence-linked assistance
      & Sources visible; user retains verification duty
      & I8, I9, I12 \\
    \addlinespace
    3 & Governed decision support
      & Context-bound evidence, verification controls, audit trail, authorized
        human reliance
      & I8--I19 \\
    \addlinespace
    4 & Constrained execution
      & Deterministic guardrails, approvals, reversibility, monitoring, full
        records
      & I8--I19 plus action controls \\
    \addlinespace
    5 & Autonomous consequential action
      & Narrowly specified settings, stringent formal controls, continuing
        assurance, explicit authority to act
      & all, plus independent assurance \\
    \bottomrule
  \end{tabular}
  \caption{Reliance tiers. The tier is set by the consequence of reliance, not by
    model quality. Tier~5 is included for completeness and should be approached
    with considerable caution in regulated settings; nothing in this paper argues
    that current systems are ready for it.}
  \label{tab:tiers}
\end{table}

Assignment to a tier needs criteria, and while a full consequence taxonomy is out of
scope, the dimensions are nameable: \emph{reversibility} of the action the assertion
supports; \emph{magnitude} of the exposure if it is wrong; \emph{breadth}, meaning how
many parties are affected and whether they consented; \emph{detectability}, whether an
error would surface before harm; and \emph{availability of a competent human check} in
the time the decision allows. A use case scoring badly on reversibility and
detectability belongs at a higher tier than its nominal risk rating suggests, because
those two dimensions determine whether anything can be done after the fact.

The principle the table exists to enforce is a negative one. Permissible reliance
does not rise because the model got better. It rises when evidence, verification,
constraints, authorization, and accountability rise together and in proportion to
what is at stake. A system that improves its benchmark score has not thereby earned
a higher tier, and a system that has implemented I8--I19 has not thereby earned
Tier~5.

Two boundaries are worth marking explicitly, because they are where programmes
drift. The step from Tier~2 to Tier~3 is the step at which the institution rather
than the user assumes the verification duty, and it requires everything in
Table~\ref{tab:invariants-system}; systems are frequently described as Tier~3 while
implementing Tier~2. And the step from Tier~3 to Tier~4 is the step from assertion
to action, at which reversibility and blast radius become first-order design
questions rather than operational details.


\section{Three Teams, Three Budgets}
\label{sec:three-teams}

Everything above is stated as though an organisation were a single agent that could
decide to satisfy Table~\ref{tab:invariants-system}. It is not, and the way it is
divided is itself a principal cause of the failures this paper describes. This
section is the least formal in the paper and, for practitioners, possibly the most
useful.

\subsection{The division}

Responsibility for the conditions of warrant is distributed across organisational
functions that were not designed to hold it jointly. In financial services the fourth
row below is a second-line function with an independent reporting line and a
supervisory mandate; in other sectors the equivalent may sit in quality assurance,
clinical governance, or internal audit. Wherever it sits, it is the function that
performs the independent challenge this paper argues for, and it is the one most often
absent from discussions of these systems.

\begin{table}[ht]
  \centering
  \small
  \begin{tabular}{@{}>{\raggedright\arraybackslash}p{2.8cm}>{\raggedright\arraybackslash}p{4.6cm}>{\raggedright\arraybackslash}p{2.4cm}>{\raggedright\arraybackslash}p{4.0cm}@{}}
    \toprule
    Function & Typical responsibility & Owns & Failure if left in isolation \\
    \midrule
    Data / knowledge
      & Source ingestion, curation, lineage, metadata, quality
      & I8, I9
      & Authentic documents with no policy applicability or output control \\
    \addlinespace
    Policy / risk / compliance
      & Rules, jurisdiction, decision authority, approvals, procedures
      & I10, I15
      & Governing policy with no control over how assertions are retrieved,
        generated, or verified \\
    \addlinespace
    Platform / AI engineering
      & Models, prompts, orchestration, retrieval, evaluation, monitoring
      & I11--I14, I17
      & Answer quality, latency, and usability with no evidentiary authority or
        institutional authorization \\
    \addlinespace
    Model risk / independent validation
      & Independent challenge, validation, and approval of models and their use
      & I13, I15, I16
      & Validating $M$ while the risk sits in $S$: a well-documented model, an
        unexamined evidence path \\
    \bottomrule
  \end{tabular}
  \caption{Four functions, and the invariants each is positioned to supply. No
    function owns a complete path from source to authorized reliance.}
  \label{tab:three-teams}
\end{table}

These are not merely three technical workstreams. They typically have separate
budgets, separate incentives, separate release cadences, separate accountability
structures, and separate systems of record. A data team is measured on lineage
coverage and freshness. A compliance function is measured on policy completeness
and audit findings. A platform team is measured on answer quality, latency,
adoption, and cost. None is measured on the grounding rate of
\eqref{eq:rate}, because it spans all three.

\subsection{The unit of failure is the handoff}

The consequence is that failures concentrate at the boundaries rather than inside
any function's remit. A system may have:

\begin{itemize}
  \item excellent data lineage and weak policy binding---I8 and I9 without I10;
  \item well-specified governing policy and uncontrolled retrieval---I10 without
        I11;
  \item high benchmark accuracy and no per-assertion provenance---nothing in
        Table~\ref{tab:invariants-system} at all, with a validation report that
        looks complete;
  \item citations without applicability validation---I12 without I11, which is
        precisely cell (v);
  \item human review without a reproducible evidence path---review that cannot be
        re-performed and therefore cannot be audited, which is I13 without I14;
  \item a validated model inside an unvalidated system---documentation of $M$'s
        training, evaluation, and limitations, with no examination of context
        resolution, applicability logic, defeater checking, or authorization. This is
        the characteristic failure of the fourth row, and it is the most dangerous
        because it produces a formal approval.
\end{itemize}

Each of these is a system in which every function has done its job. The ungoverned
handoff is the operational unit of failure, and it is invisible from inside any one
function's metrics.

This also explains a pattern worth naming, because it is the most expensive form of
the problem. Cell (v)---misapplied grounding---sits exactly on the boundary between
the data function, which supplies authentic versioned sources, and the policy
function, which knows which sources govern which decisions. The metadata required to
enforce applicability is a policy artefact stored in a data system and consumed by a
platform component. It is nobody's deliverable, and it is the single most common
thing I find missing.

\subsection{What follows}

Two conclusions, and the first is uncomfortable for how these programmes are
usually organised.

Evidentiary controls must be architecture and cross-functional invariants, not
after-the-fact review tasks. A review step owned by one function cannot establish a
property that requires artefacts from two others, and adding reviewers does not
change this. Table~\ref{tab:invariants-system} is a specification for a system, and
a specification that spans three budgets needs an owner who spans them.

A fifth function appears once retention is considered. Privacy and records
management own the tension described in Section~\ref{sec:retention}: the audit record
that reconstructability requires is in direct conflict with minimisation and erasure
obligations, and neither the data, policy, platform, nor validation function is
positioned to resolve it. That handoff has the same structure as the others and is
usually discovered later.

And the grounding rate is the right cross-functional metric precisely because no
single function can move it alone. Accuracy can be improved by the platform team
without reference to the others, which is part of why it is the number that gets
reported. \eqref{eq:rate} cannot: improving it requires corpus work, policy
metadata, and pipeline control together. A metric that only improves when three
functions cooperate is an unusually honest instrument for a problem whose failures
live between them.


\section{An Evaluation Protocol}
\label{sec:evaluation}

The preceding sections argued for a property and described an architecture that makes it
observable. This section is the operational residue: what to run, in what order, and what
the numbers mean.

Four measurements, in the order I would add them to an existing evaluation
pipeline. None replaces accuracy; all of them answer questions accuracy cannot.

\subsection{Attribution audit: the ungrounded assertion rate}
\label{sec:measure-rate}

\textbf{What it measures.} Equation~\eqref{eq:rate}, decomposed by the three rows
of Table~\ref{tab:remedies}.

\textbf{Procedure.} Sample assertions from deployment-representative traffic---not
from a curated question set, since the point is to measure behaviour on the
distribution you actually serve. For each sampled assertion:

\begin{enumerate}
  \item Identify the evidence the system had available, from the execution trace.
  \item \textbf{(G1)} Determine whether any of it supports the assertion and is
        applicable at $c$: authentic, current at $t$, right jurisdiction and
        intended use. Failure here with support present is the misapplied rate.
  \item \textbf{(G2)} Perturb: remove the putatively supporting evidence, or
        replace it with evidence for a contrary claim, and re-run. If the
        assertion survives unchanged, it was not resting on that evidence.
        Failure here is the ungrounded rate.
  \item Where the assertion is grounded and false, record it against the
        \emph{corpus}, not the system.
\end{enumerate}

\textbf{Honest limits.} Step 2 requires judgement about applicability, which means
human raters or a well-validated entailment model, and agreement on attribution is
imperfect. Report the rate with its inter-rater agreement, not as a point estimate.
And per Section~\ref{sec:measurability}, step 3 works for retrieval-provided
evidence and does not work for parametric evidence: say which you measured.

\subsection{Abstention items: catching the lucky cell}
\label{sec:abstention}

This is the single most useful addition to an existing suite, and the one that
most directly changes what the suite selects for.

\textbf{What it measures.} Whether the system asserts in the absence of evidence.

\textbf{Procedure.} Construct items whose answers are \emph{absent from the
accessible evidence}, where the only correct behaviour is abstention or an explicit
statement that the basis has not been established. Then---and this is the
part that inverts normal practice---\textbf{penalise correct answers on these
items}. A system that answers such an item correctly has demonstrated only that
its guessing distribution happened to align with the world, which is cell (ii),
and rewarding it selects for exactly the decoupling of assertion from evidence
that Section~\ref{sec:architecture} identified as the root defect.

\textbf{Construction.} Straightforward where you control the retrieval corpus:
hold out documents and ask about their contents. Harder for parametric knowledge,
where establishing absence is the same intractable problem as before. Practical
approach: build abstention items over a controlled corpus, and treat the resulting
abstention rate as a property of the retrieval-grounded configuration rather than
of the model.

\textbf{The coverage constraint, which is not optional.} Penalising true answers on
abstention items creates an obvious incentive: abstain more. Measured alone, the
statistic is maximised by a system that never answers anything. So it must be reported
\emph{jointly with coverage}---the proportion of answerable items the system does
answer---and read as a curve rather than a point, exactly as risk and coverage are
traded off in selective prediction. A system that improves its abstention score while
its coverage falls has not improved; it has become more cautious, which may or may not
be what was wanted. The pair is the measurement. Either number alone is gameable in
opposite directions.

\textbf{Why it matters more than it looks.} Most suites contain no items whose
correct answer is silence. A system therefore never encounters a case where
answering is penalised, and no optimisation pressure toward abstention exists
anywhere in the development loop. Adding these items is the cheapest available
change to what the process selects for.

\subsection{Form-invariance audit: measuring I1 directly}
\label{sec:form-audit}

\textbf{What it measures.} How far the system is from invariant I1---a number for
the property Section~\ref{sec:canonical-form} says invalidates test suites.

\textbf{Procedure.} Build a battery of formulation pairs that are semantically
equivalent and syntactically distinct: contraposed conditionals, negation duals
(``no $A$ are $B$'' / ``all $A$ are non-$B$''), quantities in words against digits,
active against passive, reordered conjunctions, identity statements in both
directions. For each pair, measure divergence in assertion behaviour at matched
output probability. Aggregate to a single \emph{form divergence} statistic---in the terms of the NLP
evaluation literature, a measure of paraphrase robustness.

\textbf{Two cautions.} Matching output probability across paraphrases is not
trivial, since the probability is itself formulation-sensitive; state the matching
method. And control for the possibility that a paraphrase is harder for reasons of
tokenisation rather than logical form, or the statistic will measure the wrong
thing.

\textbf{The use that matters.} Run the battery across model checkpoints,
post-training stages, and inference-compute budgets. This answers a question the
field currently assumes an answer to: whether additional scale, training, or
inference-time computation moves systems toward form invariance, or merely broadens
the
set of formulations handled well. Falling divergence is genuine progress on I1.
Rising accuracy with flat divergence is broader coverage, which does not restore
the validation inference of
Section~\ref{sec:validation-consequence}. I am not aware of a systematic published
measurement of this, which is itself worth noting given how much rests on the
assumption.

\subsection{Repeatability audit: measuring I7}

\textbf{What it measures.} Whether a verified answer is a verified behaviour.

\textbf{Procedure.} Re-issue identical inputs at multiple points within a session,
across sessions, and after intervening turns of unrelated content. Report the proportion
of inputs whose assertion changes, and report it separately for the three conditions,
because they fail for different reasons: within-session drift indicates context
accumulation, cross-session variation indicates a serving or version problem, and
sensitivity to intervening unrelated turns indicates that conversational state is
influencing content it should not touch.

\textbf{Why it belongs alongside the form-invariance audit.} I1 and I7 are the two
invariants that invalidate a test suite, and they do so in complementary ways.
Form-invariance failure means a passing test does not extend to the paraphrase;
repeatability failure means it does not securely cover even the string that was tested.
Measuring the first without the second gives a misleading reassurance, because a low form
divergence over inputs that are individually unstable is not evidence of anything.

\textbf{What is acceptable.} Unlike the other three measurements, this one has a defensible
target: for a system intended to reach Tier~3, repeatability under fixed version and
pinned decoding should be effectively total, because the conditions of
Section~\ref{sec:determinism} make it achievable. A non-trivial repeatability failure at
that tier is not a tolerance to be negotiated; it is a sign that one of those conditions is
not actually in place, and the audit's real value is in finding out which.

The measurement is cheap---no raters, no reference answers---and it directly quantifies how much
a passing test case is worth.

\subsection{What number is acceptable}
\label{sec:thresholds}

A measurement regime without thresholds is not a control regime, and the objection is
fair. There is no universal threshold, and I will not invent one; but the method for
deriving a local threshold is available and it does not require new data.

The threshold is set by two things: the consequence tier of
Section~\ref{sec:tiers}, and what control sits downstream of the assertion. Those two
bound the answer from opposite directions.

Where a competent reviewer examines every assertion before it is relied upon, a high
ungrounded rate is tolerable and the metric is not a safety measure at all---it is a
\emph{workload} measure, predicting how much reviewer time the system will consume and
whether the review is a rubber stamp. A system whose ungrounded rate exceeds what
reviewers can actually check has silently converted a control into a formality, and the
grounding rate is how that is detected.

Where the assertion is relied upon directly, the tolerable rate is bounded by something
the institution already knows. Every regulated firm has an accepted error rate for the
equivalent manual process: a sampling standard for file reviews, a tolerance for
unevidenced decisions found in quality assurance, an internal audit threshold that
triggers a finding. Those numbers exist, they were negotiated with a regulator or set by
policy, and they are the right anchor---because the question a supervisor will ask is not
whether the system is good in the abstract but whether it is at least as well evidenced
as the process it replaced. Anchoring to the incumbent standard also makes the comparison
honest in the direction that usually favours the system: manual processes frequently
cannot produce the evidence trail this paper requires either.

Two cautions. Do not set a threshold on pooled accuracy and infer a grounding threshold
from it, which Corollary~\ref{cor:no-bound} forbids. And set the abstention and coverage
pair jointly, per Section~\ref{sec:abstention}, or the threshold will be met by a system
that declines to answer.

Collecting the four measurements with the two framing questions, a validation report
that supports sign-off contains, at minimum:

\begin{table}[ht]
  \centering
  \small
  \begin{tabular}{@{}>{\raggedright\arraybackslash}p{4.2cm}>{\raggedright\arraybackslash}p{4.0cm}>{\raggedright\arraybackslash}p{5.4cm}@{}}
    \toprule
    Statistic & Answers & Caveat to state \\
    \midrule
    Ungrounded assertion rate
      & how often it asserts without support
      & retrieval-grounded or parametric; rater agreement \\
    \addlinespace
    Misapplied rate
      & how often support is inapplicable at $c$
      & which components of $c$ were checked \\
    \addlinespace
    Error rate (grounded, false)
      & corpus quality
      & attributed to sources, not the system \\
    \addlinespace
    Abstention rate on held-out items
      & whether silence is available
      & corpus-controlled configuration only \\
    \addlinespace
    Form divergence
      & what a passing test is worth
      & probability-matching method \\
    \addlinespace
    Repeatability
      & whether verification persists
      & within- and cross-session \\
    \addlinespace
    Accuracy
      & performance on this sample
      & does not bound any of the above (Cor.~\ref{cor:no-bound}) \\
    \bottomrule
  \end{tabular}
  \caption{Accuracy stays on the report. It moves from the headline to a line item
    with a stated scope, which is what Corollary~\ref{cor:no-bound} licenses.}
  \label{tab:report}
\end{table}


\section{Related Work}
\label{sec:related}

This paper draws on five bodies of work and departs from each in the same way: it
treats the object of interest as institutional reliance on a particular assertion
rather than as a property of a model or a distribution.

\paragraph{Attribution and citation faithfulness.} The problem of determining
whether generated text is attributable to identified sources has been given a
careful operational treatment, and the framework of \citet{rashkin2023} is what the
audit procedure of Section~\ref{sec:measure-rate} adopts. Related work on
faithfulness against factuality \citep{maynez2020} draws essentially the
distinction between grounding and truth used here, and I claim no priority for it;
surveys of the area establish it as standard \citep{ji2023}. The departure is
that measuring attribution is necessary and not sufficient for the question this
paper asks. A citation can be present, relevant, and even authoritative while still
failing temporal validity, applicability to $c$, authority for the decision type, or
authorization for reliance. Attribution establishes I12. Warrant requires I8--I19.

The sharpest form of the departure is worth stating as a result rather than as a
difference of emphasis. Faithfulness and attribution metrics evaluate whether an
assertion is supported by a supplied source. A misapplied assertion in the sense of
Section~\ref{sec:cell-v} \emph{is} supported by its source---the entailment holds,
the citation is accurate, the document is authentic and may be authoritative---and
the source does not govern the case at issue. Such an assertion therefore scores as
faithful on every metric of this family, and the failure is invisible to them by
construction rather than by oversight. Detecting it requires conditioning on a
resolved context, which these metrics do not represent.

\paragraph{Calibration, confidence, and abstention.} A substantial literature
addresses whether models know what they know and how uncertainty can be surfaced
\citep{kadavath2022,lin2022}. The results are real and the techniques useful. The
departure is the one recorded at I5 and generalised in
Section~\ref{sec:independence}: aggregate calibration is a property of a
distribution, and reliance decisions are made about individual assertions. A
calibrated probability of correctness is not an audit trail, does not identify the
evidence relied upon, does not establish applicability to the case at hand, and does
not constitute institutional permission to rely. It answers a different and easier
question.
% TODO(cite): the selective prediction / abstention literature belongs here.

\paragraph{Retrieval-augmented generation.} Retrieval conditioning
\citep{lewis2020} is the architectural ancestor of everything in
Section~\ref{sec:governed-rag}, and it does improve factuality and supply
provenance of a kind. Section~\ref{sec:retrieval-insufficient} states the
departure precisely rather than dismissively: standard retrieval supplies
authenticity where the corpus is controlled and traceability where the event is
logged, and supplies none of the three dimensions that are relations to $c$.
Governed retrieval is not a criticism of retrieval but a specification of the
controls that make retrieved evidence capable of supporting an institutional
assertion.

\paragraph{Alignment, instruction following, and guardrails.} Preference
optimisation \citep{ouyang2022} and rule-conditioned training
\citep{bai2022} shape output distributions and improve compliance behaviour, and
Section~\ref{sec:rlhf} treats the first as part of the object of analysis rather
than as a competing proposal. The departure follows from the structure of
\eqref{eq:rlhf}: the reward is a function of the output, not of the evidential
base, so these methods cannot establish a source-specific, version-specific,
context-specific evidence path for a particular claim. They can make an assertion
better formed, and can make an ungrounded assertion better formed too.

\paragraph{Assurance cases and regulated validation.} The argument here belongs to
a tradition older than the technology: structuring a claim, the evidence for it,
the argument connecting them, the controls that were applied, and the accountable
party. Safety-case and assurance-case practice in regulated engineering, and data
governance practice around lineage, authority, and temporal validity, both supply
the shape of Table~\ref{tab:invariants-system}. The contribution of this paper is to
apply that logic to machine-generated assertions specifically---to identify which
conditions the architecture of Section~\ref{sec:architecture} makes unavailable by
default, and what has to be built to restore them.
% TODO(cite): assurance-case and safety-case literature, and a data-governance


\section{Conclusion}
\label{sec:conclusion}

The paper makes two connected claims and draws one practical conclusion.

The first claim is a redefinition. Hallucination, as ordinarily defined, is a false
or fabricated output, and that definition is wrong in both directions. It convicts a
system that faithfully reports an applicable source which happens to be
wrong---an error, whose fault lies in the evidence---and it acquits a system that
asserts something true while resting on nothing, which is the same defect as
fabrication with a luckier outcome. Defined instead as ungrounded assertion, with no
truth condition, hallucination becomes a property of the relation between an
assertion and an evidential base. That is not a terminological preference. The
truth-conditional definition can only be evaluated against reference answers, offline,
on questions already answered; the redefinition can be evaluated from the system's
own execution trace, at the point of assertion, and therefore used as a gate.

From the redefinition, what a system can assert falls into five categories, not two. Crossing
grounding against truth separates a right answer reached for a reason from a right
answer reached for none, and a faithful report of a bad source from a fabrication.
Requiring grounding to be applicable at a resolved context $c = \{j, t, f, u\}$
separates a third case from both: an assertion resting on evidence that is
authentic, versioned, and authoritative, and that does not govern the decision at
hand. These five categories reduce to three remedies with three different owners,
and a pooled hallucination rate identifies none of them.

The failures are structural. The architecture maps token sequences to
distributions over token sequences, with no stage at which two encodings of one
claim are brought to a common form. The pretraining objective contains no term that
is a function of truth, evidence, provenance, or context, and none minimised by
declining to answer. In-context learning makes the prompt the entire control
surface, and the prompt is a string. Preference optimisation fits a reward to
comparisons between outputs, so grounding is not visible to it, and the pressure
runs toward assertion and away from abstention. And there is a statistical floor: a
calibrated model must assert some rare facts incorrectly, and enterprise questions
live where facts are rare.

The second claim---and it is the step that decides what to measure---is that there
is no canonical form. Nothing constructs the map to be constant on classes of materially equivalent
requests, and nothing in the objective penalises divergence between them. So a
passing test case establishes that one string produced an acceptable output on one
occasion, and not that a materially equivalent string will, nor that the same string
will later in the same session. The equivalence classes are unbounded, so coverage
is not a remedy; the failures are silent, so nothing alerts; and repeatability fails
too, so the guarantee does not securely cover even the tested case. Semantic
normalisation is not the fix, because material equivalence is context-sensitive,
depends on domain ontology, and turns on distinctions that regulated language uses
deliberately. Accuracy is a property of a sample, and the sample does not generalise
in the direction validation needs of it. A model's latent representation is
therefore not a candidate system of record: it has no stable identity, no
inspectable normalised form, no declared dependencies, and no reproducible
derivation.

Grounding does not have this defect. Checking correctness requires knowing
the answer, so it can happen only offline, on questions nobody needed to ask.
Checking grounding requires knowing what evidence the system used and whether it
governs the case---both properties of the system's own execution, available at
runtime, per assertion, on live traffic. That is the difference between a metric and
a gate, and blocking, escalating, qualifying, and logging a reconstructable basis all
require a gate. Proposition~\ref{prop:independence} shows the two properties are not
merely different but independent: no accuracy figure constrains the rate at which a
system asserts beyond its evidence, and an intervention that raises accuracy by
encouraging confident guessing lowers that rate while improving the headline number.

The practical conclusion is a requirement on architecture rather than a
recommendation about metrics. The grounding rate is well defined for any system and
measurable only where the evidential base is explicit at inference
time---identified, versioned, scope-filtered, and resolved to a context, at claim
granularity, with verification of the support relation and an abstention path that
is a branch rather than a sentence. The nineteen invariants of
Section~\ref{sec:invariants} separate what the model will not supply from what the
system must, and the division of that work across data, policy, and platform
functions is itself a principal cause of failure: the operational unit of failure is
the ungoverned handoff, and the grounding rate is the right cross-functional metric
precisely because no single function can move it alone.

None of this is an argument against using these systems. I build them, and they
work. It is an argument about what to measure, in what order, and what a passing
test entitles an organisation to conclude. A language model may be useful, fluent,
and often accurate. But in a regulated enterprise, a proposition becomes fit for
consequential reliance only when its context is resolved, its evidence is authentic
and applicable, its support relation is verified, its production path is
reconstructable, and its use is authorised within an accountable institutional
process.


\appendix
\section{A Case Traced End to End}
\label{sec:worked}

The argument has been made in fragments. This appendix traces one case through every
condition, showing what each invariant produces as a record and where the assertion
would have been stopped. The case is constructed to exercise applicability, both time
dimensions, and the defeater search at once; it is an illustration of the failure shapes
described above rather than a report of a measured incident.

\paragraph{The request.} An operator asks: \emph{``Can we rely on the accelerated
pathway for this filing?''} A system with a free-form interface will answer. Nothing in
the sentence is ill-formed.

\paragraph{Context resolution (I10).} The request is missing every component of $c$. The
system resolves, or refuses:

\begin{center}
\small
\begin{tabular}{@{}l>{\raggedright\arraybackslash}p{8.9cm}@{}}
\toprule
$j$ & jurisdiction of the filing --- \emph{not stated}; two candidate regimes apply to
      the entity \\
$t_v$ & valid time --- the filing date, 14 months ago, \emph{not stated}; the operator
      assumed today \\
$f$ & material facts --- entity class and product type, \emph{partially stated} \\
$u$ & intended use --- a binding submission, not an internal estimate,
      \emph{not stated} \\
\bottomrule
\end{tabular}
\end{center}

Four of four components unresolved or wrong. Under I10 this request does not proceed; it
returns a resolution prompt. \textbf{Most deployments stop being defensible here}, and
nothing downstream can recover it, because a well-grounded answer to the wrong question
is indistinguishable in the record from a well-grounded answer.

\paragraph{Retrieval under applicability (I11).} With $c$ resolved, retrieval filters on
scope rather than similarity alone. Similarity search returns guidance for both
jurisdictions and for the current period; the applicability predicates admit only
instruments governing $j$ whose validity interval contains $t_v$. Two of the five
highest-similarity passages are excluded: one for jurisdiction, one for
anachronism---it postdates the filing by four months and would have supported the
answer the operator expected.

\paragraph{Temporal checks (I9).} The corpus watermark $t_k$ is today; $t_v$ is 14 months
ago. By \eqref{eq:bitemporal-anachronism} the base may contain instruments not in force
at $t_v$, which is what the previous step caught. Had $t_k$ predated $t_v$,
\eqref{eq:bitemporal-stale} would flag possible omission instead, and the correct action
would be to decline rather than to answer from a base known to be incomplete.

\paragraph{Support and sufficiency (I12, I13, I18).} The surviving passage supports
eligibility for the pathway. Claim decomposition splits the response into four atomic
claims; three link to that passage. Coverage checking against the enumerated conditions
of the rule finds the passage addresses four of five: it is silent on a filing-history
condition. $\Sufficient$ fails. The correct output is a qualified answer naming the
unaddressed condition, not an unqualified yes.

\paragraph{Defeater search (I19).} The governed contrary query returns a supervisory
statement of equal standing, in force at $t_v$, narrowing the pathway for this entity
class. It does not contradict the primary passage---it \emph{undercuts} it, by removing
the entity from scope. A search for contradictions would have found nothing.
$\NoDefeater$ fails, and the case is \textsc{Contested} in the sense of
Section~\ref{sec:contested} unless the institution's precedence order separates the two
instruments. Here it does not: both are of the same standing.

\paragraph{Gating (I17).} Three conditions have failed---applicability caught two
candidate sources, sufficiency is unmet, and a defeater exists. The gate returns an
escalation with the unresolved question stated, not an answer with caveats.

\paragraph{What the record contains (I14, I15, I16).} Resolved $c$; the corpus release
and $t_k$; every retrieved passage with source identity, version, and validity interval;
the applicability exclusions and their reasons; claim-level links; the coverage gap; the
defeater and its standing; the gate decision; and the escalation recipient.

\paragraph{What each metric would have said.} This is the point of the exercise.

\begin{center}
\small
\begin{tabular}{@{}>{\raggedright\arraybackslash}p{4.3cm}>{\raggedright\arraybackslash}p{2.5cm}>{\raggedright\arraybackslash}p{6.2cm}@{}}
\toprule
Measurement & Verdict & Why \\
\midrule
Reference-based accuracy
  & possibly correct
  & The unqualified answer may match a reference answer written without $c$ \\
\addlinespace
Faithfulness / attribution
  & clean
  & The citation is accurate and the passage does entail the claim \\
\addlinespace
Context relevance
  & high
  & Topic matches exactly \\
\addlinespace
Ungrounded assertion rate
  & \textbf{hallucination}
  & Applicability, sufficiency, and no-defeater all fail \\
\bottomrule
\end{tabular}
\end{center}

Three of the four established measurements pass a case that should not have produced an
assertion at all. That is the paper's argument in one row.


\bibliographystyle{plainnat}
\bibliography{refs}

\end{document}
